\chapter{Annular comparison and width deformation}
\label{ch:annular-comparison}

This chapter supplies the comparison used by the width argument.  There are
two different passages.  First, a small annulus compares two ramp curves and
then compares the filling areas of their projected boundary loops.  Second,
the same comparison is applied to a compact family of loops after polygonal
approximation and curve shortening.  The second passage is where the
``short loop or controlled area'' alternative used in the extinction chapter
is obtained.

The carrier is retained throughout.  Let $F$ be a Ricci flow on
$[a,b]$, let $g(t)$ be its metric, and let $M$ be compact, Hausdorff and
second countable.  In the flow part the dimension is $n$ until the disk
comparison, where it is exactly three.  We use
\[
 \partial_tg=-2\operatorname{Ric},\qquad R=\operatorname{tr}_g\operatorname{Ric},
 \qquad A=|\operatorname{Rm}|.
\]
For $\rho>0$, write $S^1_\rho=\mathbb R/\rho\mathbb Z$ with its flat
metric and $P_\rho=M\times S^1_\rho$ with metric $g(t)+d\theta^2$.
All projections, ramps and annuli below use this same product.  A ramp is
an immersed periodic curve whose tangent has positive component in the
oriented circle direction.  Its circle degree is the positive integer
measuring the winding of that projection.  The canonical ramp of a loop
$\gamma$ is $x\mapsto(\gamma(x),[\rho x/(2\pi)])$, of degree one.
For an immersed curve with unit tangent $T$, write $H=\nabla_TT$ for
its curvature vector and $\Theta=\int|H|\,ds$ for its total curvature.
The ramp chapter supplies existence on the whole slab, continuous dependence
on initial curves, and constants $K_0,K_1,K_2\ge0$ controlling respectively
unit-input curvature, its covariant derivative, and Ricci curvature.

We use $\Lambda^1M=C^1(S^1,M)$ with its $C^1$ topology, where
$S^1=\mathbb R/(2\pi\mathbb Z)$, and $S^2=\{z\in\mathbb R^3:|z|=1\}$.
A family is pointwise null if each of its loops is null-homotopic; a
continuous choice of null-homotopies is not required.  For such a loop,
$A_g(\gamma)$ denotes the infimum of the Jacobian areas of Lipschitz disks
whose boundary is $\gamma$ up to a circle homeomorphism.  Existence of a disk,
nonnegativity, and continuity of this infimum in the $C^1$ topology are the
filling-area inputs used below.  Family homotopies take values in
$\Lambda^1M$ and are free, without a prescribed base loop.

\section{Annuli and their area infimum}
\label{sec:annulus-definitions}

Write $\mathbb A=[0,2\pi]\times[0,1]$ and identify its first coordinate
periodically.  We integrate on the closed rectangle
\[
 D_{\rm ann}=\{p\in\mathbb R^2:0\le p_0\le2\pi,\ 0\le p_1\le1\}
\]
which is compact and measurable.  If
$c_0,c_1:\mathbb R\to M$ are the two boundary maps, an admissible annulus
for a metric $g$ is a map $u:\mathbb A\to M$ with the following data:
\begin{enumerate}
\item $u$ is continuous on $D_{\rm ann}$ and periodic in the first
  coordinate;
\item $u(x,0)=c_0(x)$ and $u(x,1)=c_1(x)$;
\item $u$ is Lipschitz on the rectangle, with a displayed nonnegative
  Lipschitz constant;
\item $u$ is manifold-differentiable almost everywhere on the rectangle and
  its area density is integrable.
\end{enumerate}
The area is the actual integral
\[
 \operatorname{Area}_g(u)=\int_{D_{\rm ann}}
       \sqrt{\det\bigl(\langle du(e_i),du(e_j)\rangle_g\bigr)}\,dp,
\]
where the determinant is taken in the two Euclidean coordinate directions.
\lean{PoincareMT.M64Annulus}
\lean{PoincareMT.m64AnnulusArea}
\source{MT2007}\dcref{Lemma 19.15, pp. 447--449}

The admissible class is deliberately larger than the piecewise-smooth class
used to build a particular competitor.  Its least area is the guarded
infimum
\[
 \mu_g(c_0,c_1)=\inf\{\operatorname{Area}_g(u):u\text{ is an admissible
 annulus from }c_0\text{ to }c_1\}.
\]
It is not assumed to be attained.  If one annulus $A$ is supplied, then the
class is nonempty, every area is nonnegative, and
\begin{equation}
 0\le\mu_g(c_0,c_1)\le\operatorname{Area}_g(A).
 \label{eq:guarded-infimum}
\end{equation}
The first inequality is an order argument: the area density is a square root
and hence nonnegative; the second is the defining infimum inequality.  This
simple guard is used every time an infimum is compared with a future
competitor.  In particular, no minimizing annulus is silently introduced.
\lean{PoincareMT.m64AnnulusAreaRange_nonempty}
\lean{PoincareMT.m64LeastAnnulusArea_nonneg}
\lean{PoincareMT.m64LeastAnnulusArea_le_annulus}

For the interpolation competitors we additionally record a finite partition
into angular strips on each of which $u$ is $C^1$, and require each
transverse slice to be the selected minimizing geodesic between its boundary
points.  These are properties of a supplied competitor, not restrictions
on the infimum class.  A map factoring through one real
coordinate has zero two-dimensional Jacobian.  This rank-one observation is
the reason that collars used to repair a boundary trace can be made with
arbitrarily small area.
\lean{PoincareMT.m60AreaDensity_eq_zero_of_factor}
\lean{PoincareMT.M64PiecewiseC1Annulus}
\lean{PoincareMT.M64GeodesicAnnulus}

\section{Projection and evolving annular area}
\label{sec:annular-evolution}

For a product point $q=(x,\theta)$, the product metric is the orthogonal sum
of the base metric and the circle metric.  Consequently the base projection
has no larger differential norm.  If $A$ is an annulus in the product at
time $t$, its literal projection
\[
 u_{\rm proj}(z)=\operatorname{pr}_1(A(z))
\]
is an admissible annulus in $(M,g(t))$ and
\begin{equation}
 0\le \operatorname{Area}_{g(t)}(u_{\rm proj})
 \le \operatorname{Area}_{g(t)+d\theta^2}(A).
 \label{eq:projection-area}
\end{equation}
The pointwise proof compares the two-by-two Gram determinants: the first
base column is the projection of the product column and the omitted circle
component contributes a nonnegative square.  Integrating the pointwise
inequality gives (\ref{eq:projection-area}); Lipschitzness and integrability
are inherited from the product annulus.
\lean{PoincareMT.m64CircleProduct_projected_density_le}
\lean{PoincareMT.m64ProjectedDensity_integrable}
\lean{PoincareMT.m64ProjectedAnnulus_of_annulus}
\lean{PoincareMT.M64AnnulusProjection}

\begin{proposition}[Annular evolution]
\label{prop:annular-evolution}
Suppose $n\ge3$ and $c_0,c_1$ are periodic $C^2$ shrinking curves on the
closed slab, initially positive-slope ramps of degree one in the same product.
Assume an actual annulus between their initial curves is given.  This last
condition cannot be inferred from their equal circle degrees alone.
Define
\[
 \mu(t)=\mu_{g(t)+d\theta^2}(c_0(\cdot,t),c_1(\cdot,t)).
\]
The area class is nonempty at every time; $\mu$ is nonnegative and continuous
on $[a,b]$.  At every $t\in[a,b)$ it satisfies
\begin{equation}
 D^+\mu(t)\le ((2n-1)A_{\sup}(t))\,\mu(t),
 \qquad A_{\sup}(t)=\sup_{x\in M}|\operatorname{Rm}_{g(t)}(x)|.
 \label{eq:annulus-forward}
\end{equation}
Here $D^+f(t)\le c$ means that, for every $\eta>0$, all sufficiently
small $h>0$ satisfy $(f(t+h)-f(t))/h\le c+\eta$.
Separately, the unit-input bound $K_0$ gives, for $a\le s\le t\le b$,
\begin{equation}
 \mu(t)\le\exp\bigl(((2n-1)K_0)(t-s)\bigr)\mu(s).
 \label{eq:annulus-exponential}
\end{equation}
The distinction matters in the later scalar-clock comparison.
\lean{PoincareMT.M64AnnulusFlowConclusion}
\lean{PoincareMT.m64AnnulusEvolution_of_M63}
\end{proposition}

We explain the geometric variation and the removal of intersecting
boundaries.  Fix an auxiliary flat circle of circumference $1$, and place
$c_0,c_1$ in the sections of auxiliary heights $0,\varepsilon$, where
$0<\varepsilon<1$.  They are now disjoint.  Denote their annular infimum
by $\mu_\varepsilon(t)$.  A radial lift of an actual original annulus gives
a competitor in this separated class.  The nonzero degree in the original
circle excludes a compression of the annular core to a disk: a null-homotopy
would project to a null-homotopy of a nonzero-degree circle map.
The two energy bounds below give the quantitative nondegeneration
needed to construct the minimum.

The analytic construction minimizes energy over the modulus and boundary
labels together.  We describe it because the annular topology alone does
not identify the minimum with the required area infimum.  Embed the compact
target smoothly in a Euclidean space, retaining the two circle coordinates
as smooth observations.  On the fixed rectangle use the energy
\[
 \mathcal E_m(f)=\tfrac12\int
       \bigl(m|\partial_xf|_g^2+m^{-1}|\partial_yf|_g^2\bigr),\qquad m>0.
\]
The weak class consists of $W^{1,2}$ maps into the target with periodic
seam, monotone boundary labels $\sigma_i$ satisfying
$\sigma_i(x+2\pi)=\sigma_i(x)+2\pi$, and a weak real lift of the original
circle phase with its prescribed nonzero period.  Boundary labels are
allowed to be merely monotone at this stage.  Normalize their value at
$0$ into one period.  The initial annulus supplies a member of this class.
More precisely, let $H_i$ be the strictly increasing real lift of
the original circle coordinate of $c_i$.  The real weak phase has
lower trace $H_0\circ\sigma_0$ and upper trace
$H_1\circ\sigma_1+d$, where $d$ is an integral multiple of the
original circle circumference.  Its radial Green identity has
these literal traces; its angular Green identity has the prescribed
affine period as its boundary term.  The offset is allowed to
vary between competitors.  Normalizing both labels therefore
does not discard annuli with a different relative phase winding.

Two coercive estimates prevent the conformal modulus from degenerating.
On almost every horizontal slice, Cauchy--Schwarz bounds the squared
phase period by $2\pi$ times its angular Dirichlet integral.
On almost every vertical slice, choose a real linear observation of the
auxiliary circle taking distinct constant values on its two boundary
sections.  Its squared boundary difference is bounded by the vertical
Dirichlet integral.  Bounded differentials of the observations and uniform
equivalence of the ambient and intrinsic norms give constants
$\alpha,\beta>0$ with
\[
 \mathcal E_m(f)\ge\alpha m,\qquad
 \mathcal E_m(f)\ge\beta/m.
\]
Thus every competitor below an energy cap $C$ has modulus in a compact
positive interval, with strict endpoint slack obtained by using $C+1$.
The constants may depend on the chosen separation; no uniformity as
$\varepsilon\to0$ is needed for this existence step.

On that modulus interval, a minimizing sequence has bounded weak derivative
columns.  Rellich compactness gives strong convergence of the maps after
embedding, weak convergence of the columns, and almost-everywhere target
convergence; closedness of the embedded target keeps the limit in it.
Helly selection of the two normalized monotone labels and the trace
identities keep the weak limit in the same free phase class.  The real
phase period is retained, even before continuity of the labels is known.
Lower semicontinuity of the positive quadratic energy, including convergence
of the modulus and metric coefficients, gives an attained weak minimum.
Any competitor above the original cap has larger energy than the seed,
so this is a minimum over all positive moduli.

Here is the compactness control on the relative offset.  Put
$b_0=H_0\circ\sigma_0$ and $b_1=H_1\circ\sigma_1+d$,
and denote the weak vertical phase derivative by $v$.
Testing its full Green identity with the constant function one gives
\[
 \int_{D_{\rm ann}}v
       =\int_0^{2\pi}b_1-\int_0^{2\pi}b_0.
\]
Monotonicity and the affine phase period $\rho$ imply
\[
 b_i(0)\leq\frac1{2\pi}\int_0^{2\pi}b_i
                   \leq b_i(0)+\rho.
\]
The lower normalization bounds $|b_0(0)|$ by a fixed $B_0$.
If $\int v^2\leq C$, integration of $2|v|\leq1+v^2$
therefore gives
\[
 |b_1(0)|\leq B_0+\rho+
                \frac{|D_{\rm ann}|+C}{4\pi}.
\]
The upper normalized label also bounds $|H_1(\sigma_1(0))|$,
so $d=b_1(0)-H_1(\sigma_1(0))$ stays in a fixed compact
interval.  Extract its convergent subsequence together with
the trace limits.  Its limit remains in $\rho\mathbb Z$
because that subgroup is closed.  Thus the same weak class,
including its literal upper offset, survives the limiting process.

Finite real-phase energy also rules out jumps in the raw labels.
On semicircular shells about a boundary point, the trace Green
identity and Cauchy--Schwarz bound the square of a monotone phase
gap by $C E/N$ across a window of radius $\rho e^{-N}$.
A jump would give a fixed positive gap for all $N$, a
contradiction.  At the angular cut use the affine seam identity
to join the two half-windows; its boundary term is exactly the
phase period, so it introduces no extra jump.  Subtracting the
affine part then transports this continuity to every cut.
The inverse of $H_i$ is continuous, hence both original labels
are continuous.  For a sequence with bounded energy, the same
estimate and the positive lower slope of $H_i$ give a common
modulus of continuity for the labels.  Thus the Helly limit
has the actual continuous boundary trace, not merely values
defined almost everywhere on it.

Local replacements must preserve the phase constraint.  In a disk compactly
inside the cylinder, lift the circle component and cone the boundary data
in target coordinates, then restore the unchanged map outside the disk.
This leaves the global phase period and both boundary traces intact.
Minimality bounds the energy on a smaller disk by a constant times the
angular energy of a surrounding circle.  Integrating over the annulus of
radii $re^{-1}$ and $r$ gives, with $X=\mathcal E(B_r)$ and
$Y=\mathcal E(B_{re^{-1}})$, an inequality $Y\le L(X-Y)$.
Hence $Y\le qX$ for some $0<q<1$.  Iteration gives a bound
$\int_{B_s}|df|^2\le C s^\beta$ with $\beta>0$, uniformly for nearby
centers.  Morrey's averaging argument yields a continuous representative.
Target variations yield the weak harmonic-map equation; the
power decay and the same-map interior regularity argument make
this representative smooth.  Conformality is proved before
boundary differentiability.  For the actual weak derivative
columns put
\[
 U=m|f_x|^2-m^{-1}|f_y|^2,\qquad
 V=2m^{-1}\langle f_x,f_y\rangle.
\]
These are $L^1$ functions.  A small periodic horizontal source
variation changes the free labels monotonically and preserves
the real phase identities.  Differentiating its energy gives
\[
 \int (\eta'\rho U+\eta\rho' V)=0
\]
for every smooth periodic $\eta$ and smooth radial test $\rho$;
$\rho$ may be nonzero at either boundary.  Interior harmonicity
gives the second identity
\[
 \int (-m^2\eta'\rho V+\eta\rho' U)=0
\]
when $\rho$ is supported in $(0,1)$.  To obtain this for every
periodic angular test, first use tests vanishing at one cut,
then rotate the weak annulus by half a period and use a partition
between the two cuts.  Rotation preserves the minimum and
transports the original columns, so both identities concern
the same $U,V$.  Independent variation of $m>0$ gives
$\int U=0$.

For completeness, these identities force pointwise conformality
in the interior.  Pair them with $e^{ijx}$ and write $u_j,v_j$
for the resulting radial moments.  In the weak sense
\[
 v_j'=ij u_j,\qquad u_j'=-m^2ij v_j.
\]
The first identity, tested through both boundary values, gives
$v_j(0)=v_j(1)=0$ after recovering absolutely continuous
representatives.  Hence $v_j''=m^2j^2v_j$ and integration
against $\overline{v_j}$ gives $v_j=0$.  For $j\ne0$ this also
gives $u_j=0$.  For $j=0$, $u_0$ is constant and its integral
is zero by modulus balance.  Uniqueness of Fourier coefficients
for the almost every $L^1$ horizontal slice, followed by Fubini,
gives $U=V=0$ almost everywhere.  Interior smoothness makes
the Gram identities hold everywhere there, including branch
points.  In particular energy equals area.
\lean{PoincareMT.M64.auxiliaryCircle_free_phase_stress_eq_zero}

To extend through a boundary, take a small half-disk in the
covering strip.  In a target chart straightening the regular
boundary curve, replace the semicircle by its parameter cone.
The label on the diameter is its affine chord in the original
boundary parameter.  The real phase is replaced at the same
time, with diameter trace $H_i$ of that chord, retaining the
upper offset when appropriate.  The map and phase Green identities have
exactly the corresponding changes in radial boundary flux;
gluing therefore gives a member of the original free class.
Small angular energy keeps the cone in the chosen chart.
For large angular energy the fixed total energy gives the same
comparison after enlarging its constant.  Minimality thus gives
the half-disk energy decay, with the modulus factors bounded
by the fixed positive modulus interval.  Even reflection of
the weak fields and averaging produce the continuous map
with its literal trace.

Rescale the two source directions to make the conformality
equations unweighted.  In target coordinates straightening
the $C^2$ boundary curve, the transverse components have zero
trace and a quadratic Poisson bound.  Odd reflection, the
weighted energy regularity argument, and the conformality
equation recovering the remaining derivative give within-$C^1$
regularity, as developed for disks in
Section~\ref{sec:annular-plateau}.  These operations retain the
original weak values and columns.  Apply them at both radial
boundaries and again after a half-period rotation.  The two
completions agree on open overlaps because they are continuous
and equal almost everywhere there.  This gives one $C^1$ map
on the closed rectangle with exactly matching angular edges.
The two-cut construction includes the corner points and does
not presume differentiable boundary labels.

Conversely, every Lipschitz competitor admits free-label
energy approximations from above.  First attach zero-area
boundary collars and pass through polar coordinates to a map
near the closed planar annulus.  Relative mollification in
finitely many target charts preserves its boundary curves
and makes it $C^1$, with area increase smaller than any
prescribed positive error.  For this map $G$, choose a smooth
positive metric $Q$ majorizing the continuous pullback form
$G^*g$ and satisfying
\[
 \int\sqrt{\det Q}<\operatorname{Area}(G)+\eta.
\]
For completeness, the relative approximation changes the map
inside one target chart at a time, using a source cutoff times
the difference between a coordinate mollification and the original
map.  It preserves germs off the allowed region and points already
of class $C^1$.  Coordinate derivatives converge in $L^2$ and
the two-column Jacobian converges in $L^1$; uniform closeness
controls the metric coefficients.  A finite chart cover and
successively halved distance and area errors keep every image
inside its chosen chart and bound the total area increase.
The boundary collars lie outside the modified compact region.
No extension of the target-valued map through the planar hole
is required.

To construct $Q$, approximate the pullback bilinear form $A=G^*g$
within operator norm $\delta/2$ by a smooth form $B$, then set
$Q=(B+B^{\mathsf T})/2+(3\delta/2)I$.
Its quadratic forms satisfy $A+\delta I\le Q\le A+2\delta I$.
Continuity of the two-dimensional determinant on the compact
annulus makes its area tend to that of $A$ as $\delta\downarrow0$,
including where $dG$ has deficient rank.
We next construct the conformal cylinder for this metric, including
the boundary labels needed for admissibility.

\paragraph{Scalar construction of the cylinder.}
\label{sec:annular-uniformization}
Write $\Omega=\{z\in\mathbb R^2:1<|z|<2\}$, with metric $Q$.
Choose a smooth function $b$ with values zero and one on its two
boundary circles.  On $H^1_0(\Omega)$ the bilinear form
\[
 a(u,v)=\int_\Omega\langle\nabla u,\nabla v\rangle_Q\,dA_Q
\]
is coercive by the Poincare inequality and uniform ellipticity on
the closed annulus.  Lax--Milgram supplies the unique $w$ satisfying
$a(w,v)=-a(b,v)$ for all zero-boundary $v$.  Thus $H=b+w$ is
weakly harmonic, with the specified traces, and
$a(H+v,H+v)-a(H,H)=a(v,v)$.
Interior elliptic regularity gives its smooth representative.
Flattening each boundary circle gives a uniformly elliptic divergence
equation for the zero-trace correction with smooth forcing.
Tangential difference quotients first give $H^2$; differentiating
once more, with the localized forcing in $H^1$, gives $H^3$ on
smaller half-disks.  Two-dimensional Sobolev embedding supplies
a continuous classical differential up to the boundary.  The
zero Sobolev trace agrees with the continuous trace, so these
representatives patch with $H$ and retain its literal values zero
and one.

There are quantitative barriers at both circles.  For
$s(z)=|z|^2$, compactness gives $|\nabla s|_Q^2\ge c_0>0$
and $|\Delta_Qs|\le C_0$ on the closure.  For sufficiently large
$a>0$, both exponentials have positive Laplacian, since
\[
 \Delta_Q e^{\pm as}
 =e^{\pm as}\bigl(\pm a\Delta_Qs+a^2|\nabla s|_Q^2\bigr)>0.
\]
Normalize the increasing exponential to zero at $|z|=1$ and
one at $|z|=2$, and the decreasing exponential to one and zero.
The maximum principle compares them respectively with $H$ and
$1-H$.  Their positive boundary slopes give, for some $c>0$,
\[
 H(z)\ge c(|z|-1),\qquad 1-H(z)\ge c(2-|z|).
\]
In particular $0<H<1$ in $\Omega$.  Taking one-sided radial
derivatives shows that the differential of $H$ is nonzero at
both circles, with positive derivative in the increasing-radius
direction.

Orient the annulus by its planar orientation and define the
conjugate one-form $\omega=\iota_{\nabla_QH}dA_Q$.  It is closed
because $H$ is harmonic.  On the polar cover
\[
 p(r,t)=r(\cos(2\pi t),\sin(2\pi t)),\qquad
 1<r<2,\quad t\in\mathbb R,
\]
the Poincare lemma gives $dV=p^*\omega$.  Periodicity of this
differential makes $V(r,t+1)-V(r,t)$ a constant $P$, independent
of both variables.  It is the flux of $\omega$ around any circle.
This constant must be positive; it is not a choice of sign or
an assumption about the sought coordinates.

Here is the energy calculation with the boundary limits retained.
Let $U=H\circ p$, and put
\[
 J=U_rV_t-U_tV_r\ge0,\qquad
 B(r)=\int_0^1U(r,t)V_t(r,t)\,dt.
\]
Green's identity on each compact radial band gives
$\int_{[a,b]\times[0,1]}J=B(b)-B(a)$, so $B$ is increasing.
The total differential energy is finite.  If
$e(r)=\int_0^1|dH(p(r,t))|^2\,dt$, integrability supplies
circles approaching either boundary with
$(r-1)e(r)\to0$ or $(2-r)e(r)\to0$: a positive lower bound
would force the divergent integral of the reciprocal distance.
The fundamental theorem along radial lines and Cauchy--Schwarz
bound the squared flux errors by constants times these quantities.
Along such circles, $B(r)\to0$ at the inner boundary and
$B(r)\to P$ at the outer boundary.  Thus $0\le B(r)\le P$.
Nonconstant boundary data force positive energy on an interior
band, proving $P>0$.  Exhausting the annulus by the selected bands
and using integrability gives the exact identity
\[
 \int_{(1,2)\times(0,1)}J
   =\int_\Omega|\nabla H|_Q^2\,dA_Q=P.
\]

Consider the normalized map
$\Psi(r,t)=(U(r,t),V(r,t)/P)$ into $(0,1)\times\mathbb R$.
We prove that it is a diffeomorphism before using its inverse.
The local analytic input is the existence of isothermal coordinates
for a smooth planar metric, proved in
Section~\ref{sec:rf-harmonic-coordinates}.  In that construction, normalized
exponential coordinates first make the metric close to Euclidean;
a cutoff extension preserves its local values and radial Gauss
identity.  A curvature bound on a compact coordinate ball permits
the local harmonic-coordinate construction.  The first harmonic
coordinate has nonzero differential, and its conjugate and the
inverse function theorem give an isothermal chart.  This is a
local use of harmonic coordinates and imposes no completeness
hypothesis on $Q$.

In these charts a local conjugate pair for $H$ is holomorphic
(or its conjugate if the chart reverses orientation).  It is
nonconstant: the identity theorem propagated through overlapping
charts would otherwise make $H$ constant on the annulus.
The complex open mapping theorem therefore makes $\Psi$ locally
open, even at possible critical points.  It is also proper.
Indeed, a compact target set bounds $H$ away from zero and one,
so continuity and the boundary values bound $r$ away from both
circles.  On this compact radial band $V(r,t)-Pt$ is bounded
by its values for $0\le t\le1$.  A bound on $V/P$ then bounds
$t$.  The preimage is a closed subset of a compact rectangle.
The nonempty image of a proper open map is closed and open in
the connected target strip, hence is the whole strip.

Surjectivity and the deck identity imply that integer vertical
translates of $E=\Psi((1,2)\times(0,1))$ cover the unit target
square except for the image of the integer angular cuts.  That
exceptional image has planar measure zero, since it is a countable
union of smooth curve images.  Partitioning into unit horizontal
bands and translating gives $|E|\ge1$.  On the other hand the
area inequality and the energy identity give
\[
 |E|\le\int_{(1,2)\times(0,1)}|\det d\Psi|
       =P^{-1}\int J=1.
\]
Two distinct points in this source band with the same image
would have disjoint neighborhoods whose images overlap in a
nonempty open set.  Applying the area inequality to a small
closed ball about the first point and to its complement counts
this positive-area overlap twice, contradicting the unit bound.
Thus $\Psi$ is injective on the band.  Nor can $E$ overlap an
integer translate of itself: remove a small closed disk in such
an overlap, retaining its disjoint translated copy.  The remaining
set still covers the unit square modulo translations but has area
less than one, another contradiction.  Deck equivariance now gives
injectivity off the angular cuts.  A collision on the cuts would,
by openness, give an open overlap containing a point outside their
null image, already excluded.  Therefore $\Psi$ is a homeomorphism
on the whole open strip.

The same argument also removes its possible critical points.
In an isothermal chart the critical points of the nonconstant
local conjugate pair are isolated.  Its actual local inverse is
conformal away from the image of such a point and continuous at
that image.  The removable singularity theorem extends the inverse
smoothly there.  Differentiating the inverse identity excludes a
zero differential of the conjugate pair.  Consequently $dH$ is
nonzero throughout the interior and $\Psi$ has smooth inverse.
\lean{PoincareMT.M64Uniformization.exists_smooth_annular_cover_chart}

Let $F_0=p\circ\Psi^{-1}$.  The conjugate differential identity
gives, for a target tangent vector $(v_1,v_2)$,
\[
 |\nabla H(F_0)|_Q^2\,|dF_0(v_1,v_2)|_Q^2
       =v_1^2+P^2v_2^2.
\]
The continuous nonvanishing differential of $H$ has a positive
lower bound on the compact closed annulus.  Together with metric
ellipticity, this identity bounds $dF_0$ on the entire target strip.
The mean value inequality on that convex strip makes $F_0$
Lipschitz, so it extends to its closure.  Its potential and deck
identities extend by continuity.  The barriers place its two
boundary traces on the correct physical circles.

The differential $dV=p^*\omega$ likewise extends continuously to
both boundary lines, and $V$ has a Lipschitz extension $\overline V$.
On either circle, $dH$ vanishes on angular tangents and is positive
on increasing radial vectors.  The identity
$dH\wedge\omega=|\nabla H|_Q^2dA_Q$ then gives $\partial_t\overline V>0$.
For $r=1,2$ the functions
\[
 \phi_r(x)=\frac{2\pi}{P}\overline V(r,x/(2\pi))
\]
are $C^1$ increasing homeomorphisms of $\mathbb R$ with
$\phi_r(x+2\pi)=\phi_r(x)+2\pi$.  Their positive continuous
periodic derivatives are bounded below, so their inverses
$\sigma_{r-1}$ are $C^1$, Lipschitz and have positive derivative.
Extending the forward-inverse identity to the boundary gives the
literal parametrization
\[
 F(x,y)=F_0(y,x/(2\pi)),\qquad
 F(x,i)=p(1+i,\sigma_i(x)/(2\pi))\quad(i=0,1).
\]
The modulus is $m=2\pi/P>0$: the metric identity gives
$m|F_x|_Q^2=m^{-1}|F_y|_Q^2$ and
$\langle F_x,F_y\rangle_Q=0$.
The half-open rectangle $0\le x<2\pi$, $0<y<1$ maps
bijectively to $\Omega$, by the inverse and deck identities.
Change of variables therefore preserves the exact area integral,
including for compositions with maps of deficient rank.
\lean{PoincareMT.M64Uniformization.exists_lipschitz_free_modulus_conformal_cylinder}

Returning to the competitor $G$, since $G^*g\le Q$, its composition
satisfies, by the weighted diagonal Gram inequalities,
\[
 \mathcal E_m(G\circ F)\le
 \operatorname{Area}_Q(F)=\int\sqrt{\det Q}.
\]
This is an actual Lipschitz annulus: local derivative bounds for
$G$ on its compact annular image combine with Lipschitzness of
$F$ to give one bound on the rectangle.  Its angular edges match
and its boundary curves are exactly $c_i\circ\sigma_i$.
Subtracting integer multiples of $2\pi$ normalizes the labels
without changing these curves.

Project to the original circle and lift on the simply connected
closed rectangle, normalizing the lower trace to
$H_0\circ\sigma_0$.  The upper trace differs from
$H_1\circ\sigma_1$ by a constant $k\rho$, $k\in\mathbb Z$,
where $\rho$ is the original circle circumference: the continuous
difference takes values in its discrete period subgroup.
The same argument on the angular edges makes their phase
difference equal to the lower degree, here $\rho$.
The local circle isometry bounds each phase derivative by the
corresponding target column.  The convex-rectangle mean value
inequality, followed by closure, makes the lifted phase Lipschitz.
Its weak derivatives are its interior classical derivatives.
Integration by parts therefore gives the radial Green identity
with the two traces just obtained, and the angular identity with
the affine seam term $\rho\int_0^1\varphi(0,y)\,dy$ for periodic
tests $\varphi$.  Observing the full target retains the actual
map columns and their energy, while the projected phase supplies
the required phase observation.  This constructs a member of the
original weak class with the displayed energy bound.
Taking the two approximation errors to zero
gives $\mathcal E_m(f)\le\mu_\varepsilon$.
Collars that interpolate boundary labels factor through the boundary
curve and have rank at most one, so they restore literal boundary labels
without increasing area.  The resulting admissible annulus gives the
opposite inequality $\mu_\varepsilon\le\operatorname{Area}(f)$.
This proves equality with the required infimum.
\lean{PoincareMT.M64FreeWeakPhaseAnnulus.exists_strict_modulus_interval}
\lean{PoincareMT.M64FreeWeakPhaseAnnulus.weightedEnergy_attained_positive}
\lean{PoincareMT.M64.auxiliaryCircle_free_ramp_weak_minimum}
\lean{PoincareMT.M64.auxiliaryCircle_free_phase_column_power_growth}

Morrey's existence and regularity theory and Lemaire's boundary analysis
provide the analytic background.  In comparing with Lemaire's Theorem 1.7,
one must use the modification in Remark 5.5 and Lemma 5.2 for targets with
nonzero $\pi_2$; the present hypotheses do not impose $\pi_2(M)=0$.
The explicit phase class and two modulus bounds above are the specialized
construction used here.

The auxiliary phase of the minimum lifts to a real harmonic function on
the cylinder, since its boundary degree in the auxiliary circle is zero.
It has constant boundary values $0,c$ with $[c]=[\varepsilon]$ and
$c\ne0$.  Subtract $cy$ and integrate its harmonic equation against the
zero-boundary difference.  The resulting Dirichlet integral is zero, so
the auxiliary phase is exactly $cy$.  Its transverse derivative is
nonzero.  Conformality makes the two derivative columns perpendicular,
with squared norms in the positive ratio prescribed by $m$; therefore
neither column vanishes, even on the boundary by $C^1$ continuity.  This
removes all branch points for the separated minimum.  In particular the
boundary parameterizations have nonvanishing derivatives and can be used
in the first variation.

Let $\Sigma$ denote this minimum at an interior time.  Extend the prescribed
boundary velocity into it.  The interior normal first variation vanishes
by minimality.  With $k_g$ measured toward the surface interior,
\[
 \frac{d}{dt}\operatorname{Area}(\Sigma)
 =-\int_\Sigma\operatorname{tr}_{T\Sigma}\operatorname{Ric}
    -\int_{\partial\Sigma}k_g
 =\int_\Sigma(K_\Sigma-\operatorname{tr}_{T\Sigma}\operatorname{Ric}).
\]
The second equality is Gauss--Bonnet with Euler characteristic zero.
The Gauss equation gives $K_\Sigma\le\sec(T\Sigma)\le K$.
Both flat factors have zero Ricci curvature, so the original base dimension
gives $-\operatorname{tr}_{T\Sigma}\operatorname{Ric}\le2(n-1)K$.
The sum is $(2n-1)K$, regardless of the two extra circle directions.
One may use either $K=A_{\sup}(t)$ or $K=K_0$ in these pointwise estimates.
Extending the boundary variation through finitely many charts, and using
the $C^1$ boundary regularity, gives actual future annuli whose areas are
bounded by this first-order expression with an $o(h)$ remainder.  The
definition of the infimum then proves the forward inequality for
$\mu_\varepsilon$.

Here is this last order argument explicitly.  Suppose
an annulus $A_t$ attains the displayed infimum at time $t$ and, for every
$\eta>0$, every sufficiently small $h>0$ has a future competitor $B_{t+h}$
with
\[
 \operatorname{Area}(B_{t+h})\le\operatorname{Area}(A_t)+h(\lambda+\eta).
\]
The infimum is at most the area of $B_{t+h}$, so division by $h>0$ gives
\[
 \frac{\mu(t+h)-\mu(t)}h\le\lambda+\eta.
\]
This proves the one-sided derivative bound for the separated class.
To remove the auxiliary separation, projection gives
$\mu(t)\le\mu_\varepsilon(t)$.  Conversely, lift any original annulus $U$
by $(p,y)\mapsto(U(p,y),[\varepsilon y])$.  The added derivative column
changes its area by at most $C_U\varepsilon$, where $C_U$ is finite by
Lipschitzness.  Taking a competitor within any fixed positive error of
$\mu(t)$ proves $\mu_\varepsilon(t)\to\mu(t)$.

Sweeping each boundary over a short time interval constructs annuli of
area tending to zero with the time separation.  Gluing these to a competitor
at one time and comparing nearby metrics proves both nonemptiness and
continuity of the infima.  Integrate the separated forward inequality first:
on any interval where its coefficient is at most $K$,
$\mu_\varepsilon(t)\le e^{K(t-s)}\mu_\varepsilon(s)$.
Passing to the pointwise limit proves the same interval inequality for
$\mu$.  Continuity of $A_{\sup}$ on the compact slab allows $K$ to decrease
to $(2n-1)A_{\sup}(s)$ on shrinking right intervals.  Dividing the resulting
bound by $t-s$ and using $(e^{Kh}-1)/h\to K$ proves
(\ref{eq:annulus-forward}), including $s=a$.
The constant choice $(2n-1)K_0$ proves (\ref{eq:annulus-exponential}).
Passing a derivative inequality directly through a pointwise limit would
not justify this step.
\lean{PoincareMT.m64AnnulusForward_of_minimal_competitors}
\lean{PoincareMT.m64AnnulusForwardDerivativeBound.of_pointwise_limit}
\lean{PoincareMT.M64.auxiliaryCircle_free_ramp_immersed_minimum}
\lean{PoincareMT.M64.auxiliaryCircle_ramp_evolution}
\source{MT2007}\dcref{Lemma 19.15 and Corollary 19.16, pp. 447--449}

\section{Small-area comparison}
\label{sec:small-area}

The intrinsic comparison is the dimension-independent geometric core used by
the ramp estimate.  An intrinsic annulus is a smooth Riemannian metric on
the closed planar annulus $\{1\le |x|\le2\}$, smooth on a neighborhood
of that set; its two boundary lengths and its area are computed in this metric.
\begin{proposition}[Intrinsic small-area comparison]
\label{prop:intrinsic-annular}
Let $N$ be an intrinsic annulus, let $K\in\mathbb R$ be
a Gaussian curvature upper bound, let $r>0$, and let
$0<\delta<1/100$.  There is $\mu=\mu(\delta,r,K)>0$ such that, whenever
\[
 K_N\le K,\qquad r<L_0,
\]
and the total absolute turning of every first-boundary subarc of length at
most $r$ is less than $\delta$, then
\begin{equation}
 \operatorname{Area}(N)<\mu
 \quad\Longrightarrow\quad
 L_1\ge\tfrac34L_0.
 \label{eq:intrinsic-comparison}
\end{equation}
Here $L_i$ is the intrinsic length of boundary $i$ over one full period.
The order is important: $\delta,r,K$ are chosen first,
then $\mu$, then the annulus.  The statement permits arbitrary real
$K$; the proof uses $\max\{K,1\}$ in the comparison ODE.
\lean{PoincareMT.M64IntrinsicAnnulusComparison}
\source{MT2007}\dcref{Proposition 19.35, pp. 467--481}
\end{proposition}

Here are the uniform choices and the length accounting in the proof.  Put
\[
 \kappa=\sqrt{\max(K,1)},\quad
 q=\min\{r/2,1/(2\kappa)\},\quad \alpha=100\delta/q.
\]
The normal Jacobi comparison gives a positive radius $R_0$, depending only
on $K,\delta,\alpha$.  Set
\[
 h=\min\{q/100,R_0\},\qquad
 \mu=\min\{(1-\delta)^2h(q/10),1/(100\kappa^2)\}.
\]
In particular $K_+\mu+\delta<\pi/2$.  These choices precede $N$ and $L_0$.
Parameterize the first boundary by arclength $p$.  For bases whose absolute
geodesic curvature is at most $\alpha$, the inward normal exponential map
$e(p,t)$, stopped at the boundary or at height $h$, satisfies
\[
 e^*g_N\ge(1-\delta)^2(dp^2+dt^2).
\]
Indeed its transverse Jacobi field starts with value $1$ and derivative
bounded below by $-\alpha$; comparison with
$\cos(\kappa t)-(\alpha/\kappa)\sin(\kappa t)$ bounds it below by
$1-\delta$ on a sufficiently small uniform interval.  The radial field is
unit and orthogonal to it by the Gauss lemma.

\paragraph{Normal returns inside a specified region.}
We first work in a Jordan region $U$ whose closure lies in $N$ and
whose boundary consists of an inner-circle arc $I$ and one or two
embedded unit-speed sides of length at most $h$.  All circle arcs
in this argument use an angular lift of length less than one period.
This specifies which arc is meant even across the usual angular cut.

There are two elementary curvature obstructions.  A simple geodesic
loop contained in $\overline U$ bounds a Jordan subregion there:
the connected unbounded exterior of $U$ cannot enter its bounded
Jordan interior without crossing its boundary.  The same observation
gives nesting for a return joined to a subarc of $I$.  At the join
of a simple geodesic loop the exterior turning is at most $\pi$;
Gauss--Bonnet therefore gives curvature integral at least $\pi$,
contrary to $K_+\operatorname{Area}(N)<\pi$.
For an embedded normal return to $I$ with a transverse terminal
join, the initial exterior turning is at most $\pi/2$ and the
terminal turning is at most $\pi$.  The geodesic side contributes
zero turning, whence
\[
 \frac{\pi}{2}\le
 \int_{U'}K_N\,dA+\int_{I'}|k_g|\,ds.
\]
If the child base $I'$ has length at most $r$, the right side is
strictly less than $K_+\mu+\delta\le\pi/2$.
The formula can be obtained by triangulating the actual region
with both joins retained as vertices: internal edge turnings cancel,
the disk Euler count supplies $2\pi$, and the two remaining
corner defects have sum at most $3\pi/2$.
Thus every such transverse return has base length greater than $r$.

\paragraph{A confined radial family and its focusing estimate.}
\label{sec:annular-confined-focusing}
Consider a two-side collision region with base length $L\le r$,
normal side lengths $A,B\le h$, and vertex $v$.  Put
$R=L/2+\max(A,B)$, and suppose $\kappa R\le\pi/4$.
We will prove
\begin{equation}
 \cos(\kappa R)L\le
 \frac{\sin(\kappa R)}{\kappa}\int_I|k_g|\,ds.
 \label{eq:annular-confined-focusing}
\end{equation}
The construction must take place inside this region: the annulus
need not be a normal neighborhood of $v$.

First fix an interior point of the base.  Following the shorter
base subarc and then its normal side gives a path to $v$ of length
at most $R$ in $\overline U$.  Minimize length among paths with
these endpoints and this confinement.  Arclength parametrization,
compactness of $\overline U$, uniform equicontinuity and lower
semicontinuity of length give a minimizing path.  Removing a loop
would shorten it, so its arclength representative is embedded;
each of its subarcs is also minimizing under the same constraint.

Its open path avoids the two geodesic sides.  Near a smooth point
of such a side, a sufficiently small geodesically convex
neighborhood on the occupied side gives a shorter replacement
unless the path follows the side.  Equality forces agreement
with that geodesic by uniqueness, and continuation of the contact
would reach a base corner or the initial point.  The initial
point is not on a side.  At a base corner the region occupies
the quadrant between the circle and its inward normal; a path
through that corner admits a strictly shorter chord in the
quadrant.  Thus neither alternative is possible.

There is also no further contact with the base.  Otherwise take
the last such contact.  The remaining open tail avoids the whole
boundary and is a smooth unit geodesic.  Compact continuation of
its equation gives its actual endpoint derivatives.  At the last
contact its derivative is tangent to the base: in a boundary
chart, a transverse departure at an interior time of a constrained
minimizer admits a local shortening across the join.  The
collision vertex is a convex corner.  Indeed equal terminal
tangents would identify the two sides, and opposite terminal
tangents would join them to a forbidden short normal return.
For independent tangents, Gauss--Bonnet with the two right base
corners gives the occupied vertex angle as the curvature integral
plus signed base turning.  It is less than $\pi/2$ under our
budget.  The interior tail therefore arrives strictly between
the two sides; a tangent arrival would coincide with a side.
Choose the base subarc in the direction of its tangential
departure and the corresponding normal side.  Together with
the tail these bound a nested Jordan region.  Its corner
turnings have sum at most $3\pi/2$, whereas its base turning
and curvature integral sum to less than $\pi/2$.
Gauss--Bonnet is impossible.  This proves that the original
minimizer has open image in $U$, and hence is a smooth unit
geodesic up to both endpoints.  Reversing it gives a confined
radial segment from $v$ to every interior base point; at the
two endpoints use the reversed original normal sides.

We next need uniqueness in a larger class: all radial vectors
of norm at most $R$ whose whole segments stay in $\overline U$.
Choose an orthonormal frame at $v$ and write $E(z)=\exp_v(z)$
where the geodesic exists.  A smooth cutoff continuation of the
geodesic equation supplies a map on the closed vector ball;
on every confined segment it agrees with the original equation.
All curvature comparisons below concern these original segments.
The scalar Jacobi equation $j''+K_Nj=0$, $j(0)=0$, $j'(0)=1$,
compared with $\sin(\kappa t)/\kappa$, gives no conjugate point
before $\pi/\kappa$.  Thus $E$ has a local inverse at every
vector in this confined class.  Each individual segment is
embedded: a first self-return gives a nested simple geodesic
loop with curvature integral at least $\pi$.  A return to the
vertex is treated in its convex corner.  Both contradict the
curvature budget.

Suppose two confined vectors represent the same point on one
fixed radial segment.  Among noncanonical contacts choose the
earliest point on that fixed segment.  This minimum exists:
the vector ball and the confinement condition are compact,
and local inverses exclude a sequence of noncanonical vectors
converging to the canonical one.  The minimum is positive,
and the two initial segments meet only at their endpoints.
They bound a nested geodesic digon.  Gauss--Bonnet gives the
sum of its two occupied corner angles as its curvature
integral, hence less than $\pi/2$.  In particular both corners
are convex and the terminal angle is acute.

Here is the perturbation that contradicts earliest contact.
Let the two terminal vectors be $z$ and $s\theta$, where
$|\theta|=1$.  There is one inverse chart for $E$ around the
whole compact segment $[0,z]$: individual regularity and
injectivity on that segment permit shrinking and patching
the local charts.  Straightening the two initial sides makes
$z$ and $\theta$ the positive coordinate axes.  The direction
$z+\varepsilon\theta$, for small $\varepsilon>0$, enters the
digon.  A small extension of $z$ beyond its terminal point
leaves the digon, and the same holds after this perturbation.
Take the first exit of the perturbed ray from the closed
digon.  It cannot exit through the first side, by injectivity
of the chart and the positive second coordinate.  It therefore
meets the other side at a parameter strictly between $0$ and
$s$.  The new vector still has norm at most $R$: near the
terminal point, the Gauss lemma and the acute terminal angle
make the inverse radius decrease when moving backwards on
the other side; away from that point, a compact tube about
the first radial segment already has this norm bound at
contacts.  First exit guarantees confinement of the entire
new segment, not merely its endpoint.  Its positive first
coordinate makes it noncanonical.  This contradicts the
choice of $s$ and proves uniqueness.

The unique confined vector $u(p)$ representing each base point
now depends continuously on $p$.  For any convergent sequence
of base points, compactness gives a convergent subsequence
of vectors; continuity of $E$ preserves the endpoint and
whole-segment confinement.  Uniqueness identifies its limit
with the vector at the limiting point.  Local inverse charts
then make this lift smooth along the base, with its literal
endpoint values the reversed normal sides.

Finally set $f(t)=\sin(\kappa t)/\kappa$ and push the radial
vector field $f(|z|)z/|z|$ forward by $E$ along this lift,
obtaining $W$.  The Jacobi comparison gives
$j'/j\ge\kappa\cot(\kappa t)$: the Wronskian of $j$ and
$\sin(\kappa t)/\kappa$ has nonnegative derivative until
their first zero.  Radial and transverse components, together
with the Gauss lemma, consequently give for the unit base
tangent $T$
\[
 \langle\nabla_TW,T\rangle\ge\cos(\kappa R),
 \qquad |W|\le f(R).
\]
The normal endpoint values give $\langle W,T\rangle=0$ at
both ends.  Integrating
$(\langle W,T\rangle)'=\langle\nabla_TW,T\rangle
+\langle W,\nabla_TT\rangle$ proves
\eqref{eq:annular-confined-focusing}.  Endpoint continuity
suffices by first integrating on smaller closed subintervals.
This proves the quantitative estimate before its two uses below.
\lean{PoincareMT.m64Intrinsic_exists_short_base_to_vertex_minimizer}
\lean{PoincareMT.m64Intrinsic_affine_collision_minimizer_interior}
\lean{PoincareMT.m64Intrinsic_exists_earliest_confined_contact}
\lean{PoincareMT.m64_exists_earlier_radial_digon_contact}
\lean{PoincareMT.m64Intrinsic_continuous_focusing_arc_length_le_turning}
\source{MT2007}\dcref{Claims 19.39--19.43 and Remark 19.42, pp. 470--473}

\paragraph{Regional length and descent.}
For each interior base point of $I$, stop its inward normal at the
first contact with $\partial U$, or at $h$, whichever comes first.
Inward transversality and the local separation by the boundary arc
give a positive initial interval in $U$.  Continuity gives a
measurable first-contact height $H(p)$, no larger than the annular
stopping height.  The endpoint lies in $\overline U$ and, if
$H(p)<h$, in $\partial U$.  Every stopped ray is embedded.  Indeed,
a first self-intersection extracts a simple geodesic loop in
$\overline U$; uniqueness for the geodesic equation and the initial
inward direction exclude a tangential retracing.  The preceding
curvature obstruction applies to its nested region.  This reasoning
uses the specified Jordan region, and does not assume that an
arbitrary essential loop of the annulus bounds a disk.

Here is the local length estimate which forces a return to $I$.
Choose a subinterval $J\subset I$ of length exactly $q$, and suppose
that no positive normal from $J$ returns transversely to $I$ before
its regional stopping time.  Parameter transversality gives a
full-measure set of bases for which contact with the inner circle,
where the normal map is regular, is transverse.  Set $H=0$ on its
null complement.  Consequently every positive good normal avoids
the inner circle up to and including its stopping time.

Colliding embedded prefixes based in $J$ cut out a nested
first-contact region on their ordered base subinterval.  That
subinterval has length at most $q$.  Jacobi comparison first at
radius $q/2+h$ gives length less than $51q/5000$, so half its
length plus the longer side is less than $151q/10000<q/10$.
Repeating the comparison at $\varrho=q/10$ gives
\[
 |I'|\le 2\varrho\int_{I'}|k_g|\,ds.
\]
In this local use, the shortness of the base is supplied by $J$;
no exclusion of long regional bases is being used yet.
Select disjoint collision intervals and enlarge them fourfold
in the arclength coordinate.  Their union covers all collision
intervals and has length at most
$8\varrho\int_J|k_g|<q/50$.
The bases with $|k_g|>\alpha$ have length less than $q/100$.
Outside these two sets the normal strip is injective.  The area
formula and the lower metric estimate show that its bases whose
height equals $h$ have length less than $q/10$.  Thus a measurable
set $Z\subset J$ of length greater than $87q/100$ remains, with
positive height less than $h$, injective endpoints on $\partial U$,
and a regular normal map at every endpoint.

Under the supposed absence of returns to $I$, these endpoints all
lie on the other side or sides of $U$.  Their total available
length is at most $2h$.  The metric lower bound instead implies
\[
 (1-\delta)|Z|\le 2h,
\]
which contradicts $h\le q/100$ and $|Z|>87q/100$.
This length comparison does not differentiate the measurable
height function.  On each inverse chart of the normal map,
project the smooth target side to its base coordinate; the
derivative bound is the reciprocal of $1-\delta$.  The
one-dimensional area inequality on a countable collection of
such charts gives the displayed estimate.  If there are two
sides, assign a shared endpoint to the first side and partition
the remaining bases onto the second, so overlaps are counted
only once.  We have therefore constructed an embedded transverse
normal return from every such $J$, with open prefix in $U$.

Suppose now that $U$ is bounded by one geodesic side and its base
has length $L>r$.  Since $2q\le r$, arclength selects an interval
$J$ of length $q$ whose every point is at least $q/10$ from both
ends of the parent base.  The constructed return starts at
$p\in J$ and ends at another point $w$ of the parent base.
Its embeddedness gives $p\ne w$.  Its union with the interval
between $p$ and $w$ is a Jordan curve; its bounded region is
nested in $U$ by the exterior argument above.  Whichever order
$p,w$ have, this interval omits a parent end segment of length
at least $q/10$.  Its length $L'$ thus satisfies
\[
 r<L'\le L-q/10.
\]
The first inequality is the short-return obstruction.  Consider
all such nested one-side regions inside the original region,
retaining the original normal family and constants.  Their base
lengths form a nonempty set bounded below by $r$.  Choose a member
less than $q/10$ above its infimum.  The same construction gives
a member below that infimum, a contradiction.  No passage to a
limiting Jordan region is needed.  Together with the short case,
this excludes every embedded transverse regional normal return.
\lean{PoincareMT.m64Intrinsic_no_central_regional_return}
\lean{PoincareMT.m64Intrinsic_no_regional_normal_return}

For a region with two sides and base length greater than $2q$,
choose the same central interval.  The two-side endpoint estimate
constructs a transverse normal return and its nested one-side
region.  This has just been excluded.  Hence a two-side region
has base length at most $2q$.
\lean{PoincareMT.m64Intrinsic_no_long_two_side_region}

\paragraph{Global collision control.}
Two embedded normal prefixes meeting for the first time, together
with one of their cyclic base arcs, bound such a two-side Jordan
region.  To see why the chosen region is inside the annulus,
select a pair of first-contact times once and retain those same
two sides.  With each of the complementary inner-circle arcs
they form a Jordan curve.  Each bounded region is either contained
in the closed annulus or contains the open inner disk; the unbounded
outer region excludes points outside the outer circle.  If the
first bounded region contains the inner disk, the second frontier
lies in its closure.  Connectedness and unboundedness of the first
exterior place it in the second exterior, so the second bounded
region lies in the first.  Choose an interior point of the first
circle arc, which is omitted by the second frontier.  It cannot
lie in the second bounded region, since it is on the first
frontier.  It therefore lies in the second exterior, which meets
the inner disk near that point.  Connectedness of the inner disk
then puts the whole disk in that exterior.  The second region
is the required annular one.  This explains both the confinement
needed for its curvature integral and the possible choice of
either cyclic arc.

On a region with base
length at most $2q$, the radial Jacobi comparison first applies at radius
$q+h$.  It gives base length less than $101q/5000$, and hence half its
length plus the longer normal side is less than $201q/10000<q/10$.
A second application at radius $\varrho=q/10$ gives the sharper inequality
\[
 \cos(\kappa\varrho)|I|
 \le \frac{\sin(\kappa\varrho)}{\kappa}
       \int_I |k_g|\,dp,\qquad \varrho=q/10.
\]
Here Gauss--Bonnet is applied to the actual first-contact Jordan region,
whose closure lies in $N$, so its curvature integral is bounded by
$K_+\operatorname{Area}(N)$.  The two stages are needed to justify the
small radius in this inequality; small lengths of the normal sides alone
do not bound the chosen base arc.

This collision statement does not assume that every full normal ray is
embedded.  To prove that assertion, take a good base $p$.  On a sufficiently
small neighborhood of $p$, boundary curvature is at most $2\alpha$.
Two distinct embedded prefixes from that neighborhood cannot meet: the
focusing inequality would give $|I|\le4\varrho\alpha |I|<|I|$ on
the short arc; the complementary arc is excluded by its length, since
the focused arc has length at most $2q<L_0$ and the base
neighborhood can be chosen so small that its complementary arc
has length greater than $2q$.
The metric lower bound makes the normal map a local diffeomorphism along
the stopped ray.  Here is the argument retaining its closed endpoint.
In the small base neighborhood just chosen, let $\mathcal P$ consist
of pairs $(p,t)$ for which $0<t<h$, the closed prefix up to $t$
is embedded, and its positive-time points lie strictly inside the
annulus.  Prefix separation for distinct bases and embeddedness
for each single base make $e$ injective on $\mathcal P$.
A sufficiently small signed collar $\mathcal C$ at $(p,0)$ lies
in one inverse chart.  Its positive-height part belongs to
$\mathcal P$, while its nonpositive-height image has coordinate
radius at most one.  Thus $e$ is also injective on
$\mathcal P\cup\mathcal C$.

If the stopped ray first returns at times $s<t$, it is embedded
on $[0,t)$ and $(p,t)$ lies in the closure of this injectivity
domain.  There is an inverse chart about $(p,s)$ whose source
lies inside that domain.  When $s=0$ use $\mathcal C$; when
$s>0$, choose $s<v<t$.  Compactness and local inverse charts
make embeddedness of the prefix $[0,v]$ stable under a small
change of base.  Its strict confinement is stable away from zero,
and the inward collar supplies stability near zero.  This gives
the required open neighborhood of $(p,s)$ in $\mathcal P$.
For a sequence $z_j$ in the injectivity domain tending to $(p,t)$,
the images $e(z_j)$ eventually lie in this inverse chart's target.
Injectivity identifies $z_j$ with their images under that local
inverse.  Continuity therefore gives
$(p,t)=e^{-1}_{\rm local}(e(p,t))=(p,s)$, a contradiction.
This argument covers essential self-returns without claiming that every
loop in the annulus bounds a disk.

Finally, restrict to the full-measure set of bases transverse to the first
boundary.  An endpoint on that same boundary would give a transverse
regional return, excluded by the same regional argument.  A ray from
this set stopping before height $h$ therefore ends on the other boundary.
Exceptional bases may be assigned height zero for the area calculation;
they are removed again before using the endpoint image.

Both cyclic arcs must be allowed: choosing a fixed cut on the circle does
not force the first collision region to use the non-wrapping arc.
Since $\kappa\varrho\le\pi/4$, this implies
$|I|\le2\varrho\int_I|k_g|$.  A disjoint-interval selection followed by
fourfold enlargement covers all collision
arcs by a measurable set $E$ with
\[
 |E|\le16\varrho\int_{\partial_0N}|k_g|\,dp
       <32\varrho\delta L_0/q\le3L_0/50.
\]
For the cyclic cover, represent a wrapping arc by the ordinary
interval $[b,a+2\pi]$ in two periods.  Apply the weighted interval
cover there and project its two pieces back into one period.
Periodicity preserves their integrals and overlap can only decrease
the integral of their union; the seam point has zero length.
The two-period turning integral accounts for the factor two in
the bound $16\varrho\int|k_g|$.  Either chosen arc contains its
two endpoints, so every collision between retained bases is
removed by this single measurable cover.
The total-turning estimate used here follows by partitioning the boundary
into $\lceil L_0/q\rceil$ arcs of length at most $q$; since $q<L_0$,
their total turning is less than $2\delta L_0/q$.
The same estimate gives
$|\{|k_g|>\alpha\}|<L_0/50$.

Outside these two discarded sets, individual-ray injectivity and removal
of every collision arc make $e$ injective on its variable-height strip.
The area formula and the metric lower bound give
\[
 (1-\delta)^2\int_{S}\operatorname{height}(p)\,dp
 \le\operatorname{Area}(N)<(1-\delta)^2h q/10.
\]
Thus bases in $S$ with height $h$ have length less than $q/10<L_0/10$.
The remaining set $Z$ has length greater than
$(1-1/50-3/50-1/10)L_0=41L_0/50$.
Intersect $Z$ with the full-measure transverse-contact set.  Its endpoint
map lands on the other boundary, is injective, and stretches arclength by at
least $1-\delta$.  Consequently
\[
 L_1\ge(1-\delta)|Z|>(99/100)(41/50)L_0>3L_0/4.
\]
This proves the comparison once the normal return and collision constructions
are in place; their geometric hypotheses and the cyclic choice, rather than
an assumed globally embedded collar, are essential.
\lean{PoincareMT.m64IntrinsicAnnulusComparison}
\lean{PoincareMT.m64Intrinsic_exists_cyclic_retained_strip_three_losses}
\lean{PoincareMT.m64Intrinsic_good_normal_ray_injOn}
\lean{PoincareMT.m64Intrinsic_normal_collision_wide_focusing}

For the product ramps, specialize to $n\ge3$.  Given $r>0$, the selected
geometry supplies one $\mu>0$ such that for every $\rho>0$, every included
time $t$, every pair of periodic $C^2$ positive-slope ramps, and every annulus
$A$ between them with $\operatorname{Area}(A)<\mu$, the same conclusion holds:
\begin{equation}
 r\le L_0,\quad
 \text{first-boundary subarc turning}<1/200
 \quad\Longrightarrow\quad
 \tfrac34 L_0\le L_1.
 \label{eq:ramp-comparison}
\end{equation}
The threshold is selected before circumference, time, and curves.  Positive
slope gives an immersion and the periodic degree is arbitrary; degree one is
part of the evolving statement above.  We describe the error margins needed
to apply the intrinsic theorem to a minimum that is initially only $C^1$
on its boundary.

Apply the intrinsic theorem with parameters
$\delta'=7/800$, $r'=r/4$, $K=K_0$, obtaining a threshold $\mu'$;
take $\mu=\mu'/2$.  Suppose the desired length inequality fails by a gap
$d=3L_0/4-L_1>0$.  Choose
$\epsilon=\min\{\mu'/4,d/4\}$.  Smooth the two $C^2$ ramp labels
closely enough to preserve positive slope, length to the specified
$\epsilon$ margin, and turning less than $3/400$ on arcs of length at
most $r/2$.  Strictness of the original $1/200$ bound provides this
margin.  Coordinate convolution and short connecting geodesics produce
an annulus to the original curves with arbitrarily small area.
Lift the smoothed curves into distinct constant sections of an auxiliary
circle.  A radial lift of their competitor and minimization produce a
separated conformal annulus with area less than
$\operatorname{Area}(A)+\epsilon<\mu'$.

The affine auxiliary-phase argument just proved removes branch points
on the closure.  The smoothed unlabelled boundary curves permit the
following finite boundary regularity argument.  Its smoothness
hypothesis concerns those curves, and does not strengthen the
$C^2$ input curves in Proposition~\ref{prop:annular-evolution}.

Take an arbitrary boundary point of the covering strip, rescale
to a half-disk, and put the harmonic map in target coordinates
$H$.  Its complex differential satisfies a bounded first-order
equation $\bar\partial\Phi=A\Phi$.  Straightening the regular
boundary arc in a metric frame permits reflection of this
equation across the diameter.  The Cauchy-transform construction
for the reflected field gives $D\Phi\in L^2$ locally.  Thus
the actual closed-$C^1$ map has $H^2$ coordinate functions on
a smaller half-disk.  Nonvanishing of its differential and its
boundary trace $c(\sigma(s))$ also imply $\sigma'(s)\ne0$.

Choose smooth metric-normal coordinates along $c$, writing the
map as $u=(u^0,u^\perp)$.  On the diameter $u^\perp=0$.
The tangential derivative is a nonzero multiple of $c'$, and
conformality makes the normal derivative metric-orthogonal to
$c'$.  Metric-normal coordinates therefore give
$\partial_\nu u^0=0$ there.  The coordinate harmonic equation
has the form
\[
 -\Delta u^j=F^j(u,Du)
   =\sum_{i=1}^2 B^j(u)(\partial_i u,\partial_i u),
\]
with smooth coefficient forms $B^j$.  On a compact smaller
chart, $Du,B^j,DB^j,D^2B^j$ are bounded.  Since $u\in H^2$,
the weak product rule shows $F\in H^1$: each differentiated
term is either bounded or a bounded factor times $D^2u$.
The flat boundary estimate now gives $u\in H^3$, using
Dirichlet data for $u^\perp$ and Neumann data for $u^0$.

Here the Neumann gain follows from the actual normal derivative
and its zero continuous trace.  Differentiate the weak equation
in the normal direction, use the zero-trace boundary estimate
for that derivative, and recover the remaining tangential
derivatives from the original equation.  The boundary term
vanishes because this is the same normal derivative, rather
than an independently chosen weak field.
In two source dimensions $H^1$ embeds locally in $L^4$.
Consequently the new $H^3$ bound gives $D^2u\in L^4$ on a
smaller chart.  Twice differentiating $F$ produces only bounded
terms, bounded factors times $D^3u$, and products of two
$D^2u$ factors.  All belong to $L^2$, so $F\in H^2$.
A second application of the mixed boundary estimate yields
$u\in H^4$.

After a cutoff, the half-space Sobolev estimates give continuous
first and second derivatives of this same map on the closed
half-disk.  Equality with the original classical derivatives
holds in the interior by weak derivative uniqueness and
extends to the boundary by continuity; integration on the
closed convex half-disk supplies the within derivatives.
Inverting the smooth target and affine source changes gives
closed-$C^2$ regularity of the attained annulus.  This works at
every angular parameter, including the period cut.  Overlaps
agree through their common interior, hence give one $C^2$ map
on the whole closed strip.
\lean{PoincareMT.M64.RampTransport.harmonic_regular_trace_exists_closed_c2}
\lean{PoincareMT.M64.RampTransport.annulus_strip_c2_of_boundary_charts}

Trim both boundaries into its smooth interior.  As the
trim tends to zero, their lengths and turning integrals converge to those
of the labelled boundary curves.  The lower length stays above $r/4$,
and arcs of length at most $r/4$ have turning less than $7/800$.
The induced metric on the trimmed cylinder is a smooth intrinsic annulus;
minimality and the Gauss equation bound its Gaussian curvature by $K_0$,
and its area is no greater than that of the minimum.  The intrinsic theorem
therefore applies to every sufficiently small trim.
Letting the trim vanish and restoring the smoothed labels yields
$3L_0/4\le L_1+2\epsilon$.  Since $2\epsilon\le d/2$, this contradicts
the definition of $d$.  All perturbations were chosen after the universal
threshold; none changes its dependence on $r$ and the ambient geometry.
\lean{PoincareMT.M64RampSmallAnnulusComparison}
\lean{PoincareMT.m64RampSmallAnnulusComparison_of_geometry}
\source{MT2007}\dcref{Lemma 19.31, pp. 462, 464--466}
The external analytic inputs are Morrey's ICM lecture, Sections 2--5,
printed pp. 180--187, and Lemaire's Theorem 1.7, p. 94, with Remark 5.5,
p. 102.

\section{Static approximation and finite nets}
\label{sec:static-approximation}

Let $\Gamma:S^2\to\Lambda^1M$ be a continuous family of $C^1$ loops and let
$\zeta>0$.  There is $N_0>0$ so that every $N\ge N_0$
and every $\gamma$ in a prescribed compact subset of the $C^1$ loop space has
one unique sampled minimizing geodesic $N$-gon $P_\gamma$.  Its vertices are
the samples at the fixed cell points, and
\begin{equation}
 \operatorname{Length}(P_\gamma)=
 \sum_{j=0}^{N-1}(2\pi/N)\,\operatorname{speed}(P_\gamma|_{[2\pi j/N,2\pi(j+1)/N]}).
 \label{eq:polygon-length}
\end{equation}
The loss is nonnegative and less than $\zeta$.  The geodesic interpolation
annulus from $\gamma$ to $P_\gamma$ has area in $[0,\zeta)$ and is both
piecewise $C^1$ and transverse-geodesic.  If $\Gamma$ is pointwise null, its
raw boundary maps have genuine filling disks and the filling-area error is
less than $\zeta$.  A fixed-profile flattening produces a $C^1$ family smooth in
the angular coordinate, with continuous first and second angular jets, and
preserves the family homotopy class.

We give the uniformity argument.  Compactness of $M$ gives a neighborhood
of the diagonal on which the minimizing geodesic between two points is
unique and depends smoothly on them.  Compactness of the loop family in
$C^1$ bounds all angular speeds and makes their first derivatives uniformly
continuous in finitely many coordinate charts.  For sufficiently fine
equal meshes, every sampled subarc stays in one of these geodesic
neighborhoods.  Each polygon side is the unique constant-speed minimizing
geodesic between its endpoints, including a constant side when endpoints
coincide.  Its speed integral is exactly the endpoint distance, which is
at most the length of the original subarc.  Uniform continuity of the
derivatives makes the sum of these chord deficits tend to zero uniformly
in the compact family.  Choose a common mesh cutoff for this length
estimate and for geodesic uniqueness.

Interpolate pointwise from the loop to its polygon by these short minimizing
geodesics.  On each open mesh cell the angular derivative is bounded by a
common constant, while the transverse derivative has norm equal to the
endpoint distance, which tends uniformly to zero.  The two-column area
bound $|\partial_xf\wedge\partial_yf|\le|\partial_xf|\,|\partial_yf|$
therefore makes the annular area tend to zero.  Across a mesh endpoint the
maps agree; there are only finitely many such endpoints, of measure zero,
so the cellwise Lipschitz and area estimates glue.  Taking the maximum of
the three cutoffs proves all assertions for every subsequent mesh size.

The polygon may have corners and need not itself belong to $\Lambda^1M$.
On each cell compose its geodesic side with a fixed smooth increasing
transition $\chi:[0,1]\to[0,1]$, flat to all orders at both endpoints.
This preserves the image and length by the one-dimensional substitution
formula and makes all angular derivatives agree at vertices.  Dependence
of geodesics on endpoints gives continuous first and second angular jets
over $S^2$.  Interpolating the time parameter between the identity and
the flattening would still leave polygonal corners.  Instead use the
original family $\Gamma$ and the flattened family $\widetilde\Gamma$.
They are uniformly close: the flattening stays in each original mesh
cell, where both the loop and its short polygon side have uniformly
small diameter.  Join corresponding values by the unique short
geodesic in the fixed diagonal neighborhood.  This gives
\[
 H(\tau,z)(x)=\operatorname{geo}_{\Gamma(z)(x),
             \widetilde\Gamma(z)(x)}(\tau).
\]
Each loop in this homotopy is $C^1$.  The chain rule expresses
its angular derivative as the endpoint differentials of the
geodesic interpolator applied to the two original angular
derivatives.  Values and first derivatives are jointly continuous
in $(\tau,z,x)$, hence the homotopy is continuous in the $C^1$
loop space.  It has the literal original and flattened families
as endpoints and transfers pointwise nullity.

Disk area requires a separate check from length.  The piecewise polygon has
a continuous raw circle boundary, with filling disks allowed up to boundary
homeomorphism.  The flattening is an orientation-preserving circle
homeomorphism, although its inverse need not be differentiable at vertices.
Absorbing it into the permitted boundary label identifies the raw polygon
disk-area range with that of the flattened loop.  Thus one transports actual
disks without precomposing the disk interior by a possibly nonsmooth inverse.
Uniform closeness, the sum of the two length bounds and the controlled
geodesic interpolation give a small-area annulus between the original
and flattened loops.  Gluing to actual disks in both directions gives
the filling-area error for these $C^1$ loops.  The equality of admissible
area ranges then transfers exactly that error to the raw polygon.
\lean{PoincareMT.M64StaticApproximationTheory}
\lean{PoincareMT.M64RawFamilyApproximation}
\lean{PoincareMT.M64RawFamilyApproximation.toM63}
\source{MT2007}\dcref{Definition 19.18, Claims 19.19--19.20 and Corollary 19.21, pp. 450--451; Lemma 19.17 and Claim 19.22, pp. 449--453}

When a polygonal approximation is already supplied, construct shrinking
ramps from its actual flattened curves and apply the analytic estimates
to those solutions.  The number of sides, the initial family, and the
annuli giving its area error remain those of the supplied approximation.
Choosing a second approximation here would lose the initial-area statement
needed in the deformation theorem.
\lean{PoincareMT.m64EvolvingApproximation_from_M63}

The finite-net output is stronger than pointwise compactness.  For every
$\mu>0$, a finite list $z_1,\ldots,z_m\in S^2$ and a circumference cutoff
$\rho_* >0$ are chosen so that each $z\in S^2$ has one node $z_i$ with, for
every $0<\rho<\rho_*$, an actual annulus at the initial time joining the
canonical ramp of $\Gamma(z)$ to that of $\Gamma(z_i)$ and having area less
than $\mu$.  The node is chosen before circumference and is retained for all
smaller circumferences.  This is the compactness step that later permits one
common circumference for the entire family.

To prove this stronger order of choices, bound the angular derivatives of
the compact family by $S$.  The minimizing interpolator has a uniform
endpoint differential bound $B$ on a smaller diagonal neighborhood.
If two loops are uniformly within $\epsilon$ there, their interpolation
annulus has column bounds $B,\epsilon$.  Its canonical lift to $P_\rho$
has bounds $B+1,\epsilon$ for all $0<\rho\le2\pi$, since its angular
circle speed is $\rho/(2\pi)\le1$ and its transverse circle derivative
is zero.  Its area is at most $2\pi(B+1)\epsilon$.  Choose $\epsilon>0$
so that this is less than $\mu$, while remaining in that neighborhood.
The image of $S^2$ is compact in the uniform topology, so finitely many
uniform $\epsilon$-balls with centers in the image cover it.  Take the
corresponding source parameters as nodes.  The same chosen node and the
same interpolator work at every smaller circumference; $\rho_*=2\pi$
is an available cutoff.  Null-homotopy is not needed for this net.
\lean{PoincareMT.M64FamilyAnnulusNet}
\lean{PoincareMT.m64FamilyAnnulusNets_of_compact}
\source{MT2007}\dcref{pp. 461--462}

\section{Deforming a raw family}
\label{sec:width-deformation}

For a flow on $[t_0,t_1]$ define its area-comparison profile by
\begin{equation}
 \Phi_F(A_0,t)=e^{-\int_{t_0}^{t}R_{\min}(s)/2\,ds}
 \left[A_0-2\pi\int_{t_0}^{t}
 e^{\int_{t_0}^{s}R_{\min}(v)/2\,dv}\,ds\right],
 \quad R_{\min}(s)=\inf_{x\in M}R(x,s).
 \label{eq:comparison-profile}
\end{equation}
It satisfies $\Phi_F(A_0,t_0)=A_0$ and
$\partial_t\Phi_F=-2\pi-R_{\min}(t)\Phi_F/2$ on the closed slab.
Differences of initial values are multiplied by the positive exponential
factor, so initial area errors transfer quantitatively to terminal errors.

\begin{theorem}[Deformation of a null loop family]
\label{thm:loop-family-deformation}
The input is a compact Hausdorff second-countable smooth three-manifold, a flow on
$[t_0,t_1]$ with $t_0\le t_1$, and a raw continuous family
$\Gamma:S^2\to\Lambda^1M$ of pointwise null loops.  The output for each
$\zeta>0$ is a family $\Gamma_t$ indexed by the closed interval, jointly
continuous in $(t,z)$, pointwise null, and freely homotopic to $\Gamma$ for
every $t$.  At the initial time,
\begin{equation}
 \left|A_{g(t_0)}(\Gamma_{t_0}(z))-
        A_{g(t_0)}(\Gamma(z))\right|<\zeta,
 \label{eq:initial-area-error}
\end{equation}
where $A_g$ is the filling-area infimum.  At the terminal time every member
satisfies the pointwise alternative
\begin{equation}
 \operatorname{Length}_{g(t_1)}(\Gamma_{t_1}(z))<\zeta
 \quad\text{or}\quad
 A_{g(t_1)}(\Gamma_{t_1}(z))
 \le \Phi_{F}\bigl(A_{g(t_0)}(\Gamma_{t_0}(z)),t_1\bigr)+\zeta.
 \label{eq:terminal-alternative}
\end{equation}
\lean{PoincareMT.M65RawFlowInput}
\lean{PoincareMT.M65DeformedFamily}
\lean{PoincareMT.areaComparisonProfile}
\lean{PoincareMT.areaComparisonProfile_hasDerivWithinAt}
\source{MT2007}\dcref{Definition 18.23 and Proposition 18.24, pp. 433--434; Claim 18.26, p. 434}
\end{theorem}

For its proof write $a=t_0$ and $b=t_1$ and initially assume $a<b$.
Choose the static approximation with error $\zeta$ and its ramp solution
family for every $\rho>0$.  Until a common circumference is selected,
$\Gamma_t(z)$ denotes the projection of the solution for the circumference
under consideration.  All constants below are uniform for $0<\rho<1$.

\subsection{The swept annulus and disk comparison}
\label{sec:swept-area}

For strict interior times $a<s\le t<b$, a shrinking ramp $c(x,r)$ gives the
time-reparameterized map
\[
 U(x,y)=c\bigl(x,s+(t-s)\chi(y)\bigr),
\]
where $\chi$ is a smooth monotone transition with $\chi(0)=0$ and
$\chi(1)=1$.  The time value stays in $[s,t]$, so the map never evaluates the
flow outside the chosen slab.  It is smooth on the rectangle and has exact
boundary curves at $s$ and $t$; hence it is an admissible annulus.
\lean{PoincareMT.m65SweptMap}
\lean{PoincareMT.m65SweptMap_contMDiff}
\lean{PoincareMT.m65InteriorSweptAnnulus}
\source{MT2007}\dcref{Claim 19.23, pp. 453--454}

The angular derivative is the ramp tangent and the vertical derivative is the
time derivative multiplied by $(t-s)\chi'(y)$.  The metric evolution bounds
the two column norms uniformly in circumference; integrating the resulting
Jacobian estimate gives
\begin{equation}
 \operatorname{Area}(U)\le
 e^{2K_2(b-a)}B_{\Gamma}(t-s),
 \label{eq:swept-area}
\end{equation}
where $B_{\Gamma}$ is the uniform total-curvature bound of the selected
solution family.  More precisely, if $C_{\Gamma}$ is its initial
length-plus-total-curvature bound and $C_1(K_0)$ is the curvature constant
in the ramp chapter, one may take
$B_{\Gamma}=2C_{\Gamma}\exp(|C_1(K_0)+K_2|(b-a))$.
The curvature-vector velocity gives the pointwise Jacobian bound
$|\partial_xc|\,|H|\,dt$, whose angular integral is the total curvature
$\Theta(t)$.  The factor $e^{2K_2(b-a)}$ transports its area to any fixed
metric time in the slab.  The estimate is linear in the time separation and is
uniform in the member and circumference.
\lean{PoincareMT.m65FamilyInteriorSweptAnnulus_area_le}
\lean{PoincareMT.m65FamilyTotalCurvatureBound}

Projection of $U$ supplies an annulus between the two base loops.  If $D_0$
is any actual Lipschitz filling disk for the first boundary, the projected
annulus glues to $D_0$ and yields a disk $D_1$ for the second boundary.  The
two directions are retained, and the guarded infimum argument gives
\begin{equation}
 \left|A_{g(q)}(\gamma_1)-A_{g(q)}(\gamma_0)\right|
 \le \operatorname{Area}_{g(q)}(U_{\rm proj}).
 \label{eq:disk-gluing}
\end{equation}
The disk construction initially has an arbitrarily small positive area
error, arising from matching the boundary labels on the gluing collar.
Given $\epsilon>0$, choose a source disk with area less than
$A_{g(q)}(\gamma_0)+\epsilon$ and glue with error $\epsilon$.
Then $A_{g(q)}(\gamma_1)\le A_{g(q)}(\gamma_0)+
\operatorname{Area}(U_{\rm proj})+2\epsilon$; let $\epsilon$ decrease to
zero.  Reverse the annulus for the other direction.  A displayed source
disk establishes nonemptiness of the target class before either infimum
is used.  In the family application, nullity and the static approximation
supply these disks at the initial time, and sweeping supplies them later.
\lean{PoincareMT.m65ProjectedDiskComparison}
\lean{PoincareMT.m65FillingArea_le_add_annulus}
\source{MT2007}\dcref{Claim 19.23, pp. 453--454}

\subsection{Changing the metric and including endpoints}

A disk at time $s$ can be transported through the same carrier to time $t$.
The Ricci metric comparison gives
\begin{equation}
 A_{g(t)}(\gamma)\le e^{2K_2|t-s|}A_{g(s)}(\gamma).
 \label{eq:metric-area-transport}
\end{equation}
Combining this with (\ref{eq:disk-gluing}) and the swept-area bound yields a
uniform estimate, valid for all included endpoint times,
\begin{equation}
 \left|A_{g(q)}(\Gamma_t(z))-A_{g(q)}(\Gamma_s(z))\right|
 \le e^{2K_2(b-a)}B_{\Gamma}|t-s|.
 \label{eq:time-area-lipschitz}
\end{equation}
Interior estimates are first proved on strict subintervals.  Continuity of
the filling area in the fixed metric then passes to $s=t_0$ or $t=t_1$;
this is why the endpoint claim does not assume a flow extension beyond the
closed slab.
The fixed-metric continuity can be obtained directly from the
sphere-family statement proved in the filling-width chapter.
Rescale $[a,b]$ to $[0,1]$.  The height map
$h:S^2\to[0,1]$, $h(z)=(1+z_3)/2$, has the continuous section
$s(u)=(2\sqrt{u(1-u)},0,2u-1)$.
Composing a null-loop time path with $h$ makes a null sphere
family; continuity of its filling area, restricted along $s$,
gives continuity along the original closed time path.
\lean{PoincareMT.m65FillingArea_flowMetric_le}
\lean{PoincareMT.m65ClosedTimeFillingDifference}
\lean{PoincareMT.m65UniformTimeFillingBounds}

Let $f_z(t)=A_{g(t)}(\Gamma_t(z))$.  The same estimates also give constants
$A_*,L_*$, independent of $z$ and $0<\rho<1$, such that
$0\le f_z(t)\le A_*$ and $|f_z(t)-f_z(s)|\le L_*|t-s|$.
For example, the upper area bound follows by transporting the uniformly
bounded initial areas and adding the swept area.  In comparing two times,
use $e^{2K_2d}-1\le2K_2e^{2K_2(b-a)}d$ for $0\le d\le b-a$.
Then $L_*=2K_2e^{2K_2(b-a)}A_*+e^{2K_2(b-a)}B_{\Gamma}$ suffices.
These uniform bounds, rather than a derivative bound assumed for $f_z$,
will control the time intervals on which the area comparison is unavailable.

\subsection{A minimizing disk with its actual boundary trace}
\label{sec:annular-plateau}

Fix a smooth metric $g$ on a compact smooth three-manifold $M$.
Let $\gamma:S^1\to M$ be smooth, injective, and regular, and suppose
that it has a Lipschitz spanning disk.  The boundary parameter of a
spanning disk is allowed to be any circle homeomorphism, as in
Chapter~\ref{ch:filling-width-v4}.

\begin{theorem}[Attainment for a regular Jordan loop]
\label{thm:annular-plateau}
There is a Lipschitz map $f:\overline{\mathbb D}\to M$ and a circle
homeomorphism $\beta$ such that $f|_{S^1}=\gamma\circ\beta$ and
\[
 \operatorname{Area}_g(f)=A_g(\gamma).
\]
The map is smooth, harmonic, and weakly conformal on $\mathbb D$,
is $C^1$ within $\overline{\mathbb D}$, and has only finitely many
zeros of its differential on the closed disk.  The assertion allows
zeros on the boundary.
\lean{PoincareMT.m65Plateau_attainment}
\end{theorem}

We give the construction, including the boundary properties needed for
area variation.  Choose a smooth embedding $e:M\hookrightarrow\mathbb
R^N$.  Extend the metric on each embedded tangent space to a positive
quadratic form $G_p$ on $\mathbb R^N$.  Compactness gives constants
$0<k\le K$ with
$k|v|^2\le G_p(v,v)\le K|v|^2$ for all $p,v$.
The auxiliary embedding serves to define weak derivatives; on a
genuine map its energy is the intrinsic energy.

A weak disk consists of an almost everywhere $M$-valued map $u$,
its embedded value $e\circ u\in W^{1,2}(\mathbb D,\mathbb R^N)$,
and a continuous weakly monotone map $\beta:S^1\to S^1$.
Here weakly monotone means a uniform limit of circle homeomorphisms,
with either orientation permitted.  The boundary value is linked to
the weak derivatives by the coordinate Green identities
\[
 \int_{\mathbb D} (e\circ u)_j\,\partial_i\phi
   +\int_{\mathbb D} d_i(e\circ u)_j\,\phi
 =\int_{S^1}(e\circ\gamma\circ\beta)_j\,\phi\,n_i\,d\theta.
\]
Thus the boundary function is the actual Sobolev trace.  Its energy is
\[
 E(u)=\frac12\int_{\mathbb D}\sum_{i=1}^2
                   G_{u}(d_i(e\circ u),d_i(e\circ u)).
\]
Fix distinct $a,b,c\in S^1$ and impose
$\beta(1)=a$, $\beta(-1)=b$, and either $\beta(i)=c$ or
$\beta(-i)=c$.  These three conditions prevent loss of the boundary
under conformal automorphisms; they do not fix the remaining boundary
parameter.

Take an energy-minimizing sequence in this normalized class.
Uniform ellipticity bounds its two weak derivatives in $L^2$.
Its boundary maps have a common modulus of continuity.  To see the
point of the three pins, take a small circular crosscut near a
boundary point.  Polar integration and Cauchy--Schwarz give a radius
between $\epsilon$ and $R$ for which the squared distance between the
two endpoint images is at most
\[
 \frac{2\pi B}{\log(R/\epsilon)},
 \qquad B\ge\|d_1(e\circ u)\|_2^2+
                    \|d_2(e\circ u)\|_2^2.
\]
The continuous inverse of the embedded Jordan curve turns this into
closeness of their circle labels.  A short source arc has diameter
less than the separation of the three pinned source points, so it
contains at most one pin.  Weak monotonicity and the other two pins
exclude the long target arc between its endpoint labels.  The whole
short source arc therefore has small image.  The crosscut scale is
chosen from $B$ and the Jordan inverse modulus, uniformly in the
sequence.

Arzela--Ascoli now gives uniform convergence of the boundary labels.
The weakly monotone class and both alternatives for the third pin are
closed.  Rellich compactness gives strong $L^2$ convergence of embedded
values, and Hilbert weak compactness gives weak $L^2$ convergence of
both derivatives.  A further almost everywhere subsequence has values
in the closed set $e(M)$.  Uniform boundary convergence also gives
strong convergence of the boundary values in $L^2(S^1)$, so all the
Green identities pass to the limit.

The metric energy is lower semicontinuous despite its varying
coefficients.  Write $G_{u_j}$ as the positive self-adjoint operator
$T_j$, acting by multiplication on $L^2$.  Almost everywhere value
convergence and the uniform coefficient bound imply
$T_jv\to Tv$ strongly for each fixed $v\in L^2$.  If $d_j\rightharpoonup d$,
positivity gives
\[
 \langle d_j,T_jd_j\rangle
 \ge 2\langle d_j,T_jd\rangle-\langle d,T_jd\rangle
 \longrightarrow\langle d,Td\rangle.
\]
Summing over the two columns proves that the limit is a normalized
weak minimum.
\lean{PoincareMT.m65NormalizedWeakDisks_boundary_equicontinuous}
\lean{PoincareMT.m65NormalizedWeakDisks_subsequence}
\lean{PoincareMT.m65WeakDisk_energy_le_of_limit}

The energy of this minimum is at most $A_g(\gamma)$.  This requires a
comparison with every admissible Lipschitz disk $D$, whose arbitrary
boundary homeomorphism need not be Lipschitz.  Let $H$ be its measurable
positive semidefinite Gram matrix, extended by zero outside the disk.
For every $\eta>0$ choose a smooth positive definite matrix field $Q$
such that
\[
 \int_{\mathbb D}\frac12\sqrt{\det Q}\,
               \operatorname{tr}(Q^{-1}H)
       <\int_{\mathbb D}\sqrt{\det H}+\eta.
\]
One first regularizes $H$ by a positive scalar matrix and then
smooths it, making $Q$ scalar near the center and outside the disk.
For fixed regularization the inverses are bounded and smoothing
converges in the displayed integral; letting the regularization tend
to zero gives $\sqrt{\det H}$, including at rank-deficient points.
The scalar regions can be chosen with arbitrarily small area error.
The following construction supplies the required coordinates on the
whole disk, including its boundary.

\paragraph{Smooth disk coordinates.}
\label{sec:annular-disk-beltrami}
Let $Q$ be a smooth positive definite symmetric matrix field on the
plane, equal to $\delta I$ near zero and for $|z|\ge1$, with $\delta>0$.
We construct a smooth plane diffeomorphism $\Phi$, preserving both
the open and closed unit disks, such that
\[
 (D\Phi_z)^T Q(\Phi(z))D\Phi_z=c(z)I,\qquad c(z)>0,
 \quad |z|\le1.
\]
Identify the plane with $\mathbb C$, and put
$\partial=(\partial_x-i\partial_y)/2$ and
$\bar\partial=(\partial_x+i\partial_y)/2$.  The coefficient
\[
 \mu(z)=
 \frac{Q_{11}(z)-Q_{22}(z)+2iQ_{12}(z)}
      {Q_{11}(z)+Q_{22}(z)+2\sqrt{\det Q(z)}}
\]
is smooth, has modulus less than one, and vanishes near zero and
outside the open disk.  Positivity also gives the algebraic identity
\[
 Q_z(v,v)=d(z)|v+\mu(z)\bar v|^2,\qquad
 d(z)=\frac{\operatorname{tr}Q(z)+2\sqrt{\det Q(z)}}4>0.
\]
Write $Jz=1/\bar z$ for reflection in the unit circle.  Extend the
coefficient by reflection, setting
\[
 \nu(z)=\mu(z)+(z/\bar z)^2\overline{\mu(Jz)}
 \quad(z\ne0),\qquad \nu(0)=0.
\]
The second summand vanishes inside the disk and the first outside.
Smoothness across the unit circle follows from the smooth vanishing
of $\mu$ on the outside; near zero the reflected summand is zero.
Vanishing of $\mu$ near zero makes $\nu$ compactly supported.
Thus $|\nu|\le k<1$ for some $k$, and
$\nu(z)=(z/\bar z)^2\overline{\nu(Jz)}$.

Here is the analytic construction for this smooth compactly
supported $\nu$.  The Beurling operator $S=\partial\bar\partial^{-1}$
on $L^2(\mathbb C)$ has Fourier multiplier $\bar\xi/\xi$ away
from zero.  Plancherel gives $\|S h\|_2\le\|h\|_2$.
Consequently $T h=\nu S h$ has norm at most $k$, and
\[
 h=\sum_{j=0}^{\infty}T^j a,\qquad a=\partial\nu,
 \qquad h-\nu S h=\partial\nu
\]
converges uniquely in $L^2$.  Smoothness of this particular solution
needs a derivative estimate.  Let $B_m(h)$ be the sum of the $L^2$
norms of all ordered coordinate derivatives through order $m$.
Commutation of $S$ with those derivatives and the product rule give
\[
 B_{m+1}(\nu S h)\le k B_{m+1}(h)+C_m B_m(h).
\]
All derivatives landing on $\nu$ occur in the lower order term.
Choose $r=(1+k)/2$.  Induction on $m$ gives
$B_m(T^j a)\le A_m r^j$: at the next order choose
$A_{m+1}\ge B_{m+1}(a)$ and
$(r-k)A_{m+1}\ge C_mA_m$, and then induct on $j$.
Sobolev estimates on the plane bound every uniform derivative norm
by one of these finite derivative budgets.  The series therefore
converges with all derivatives.  Every term has support in
$\operatorname{supp}\nu\cup\operatorname{supp}a$, so its sum is
smooth and compactly supported.

The Cauchy potential
\[
 q(z)=\frac1\pi\int_{\mathbb C}\frac{h(\zeta)}{z-\zeta}\,dA(\zeta)
\]
satisfies $\bar\partial q=h$, $\partial q=S h$, and $q(z)\to0$
at infinity.  These identities follow first distributionally
from the fundamental solution of $\bar\partial$ and then
classically from the smooth right side; equivalently one may
differentiate after subtracting the value at the kernel's singularity.
Compact support gives the decay directly outside a large disk.
It follows that
$\bar\partial q-\nu\partial q=\partial\nu$.
Hence $w=e^q$ never vanishes, tends to one at infinity, and satisfies
$\bar\partial w=\partial(\nu w)$.  The complex-valued real one-form
\[
 A_z(v)=w(z)(v+\nu(z)\bar v)
\]
is closed.  Its actual primitive is
$f(z)=\int_0^1 A_{tz}(z)\,dt$.
Differentiating this integral and using closedness changes its
derivative integrand to
$\frac d{dt}\{t A_{tz}(v)\}$, so $Df_z(v)=A_z(v)$.
In particular $f$ is smooth and
\[
 \bar\partial f=\nu\,\partial f,\qquad
 \partial f=w,\qquad
 \det Df=|w|^2(1-|\nu|^2)>0.
\]

This local diffeomorphism is global.  Indeed $Df\to I$ at infinity.
Choose $R$ so that $\|Df-I\|\le1/2$ outside $B_R$ and integrate
along radial segments starting on its boundary.  A uniform bound
there gives $|f(z)-z|\le |z|/2+C$, and hence
$|f(z)|\ge |z|/2-C$.  Thus $f$ is proper.  Its image is nonempty,
open and closed in the plane, and its fibers are finite.
Local inverse neighborhoods and properness make it a covering:
after taking the finitely many inverse neighborhoods over one
point, properness excludes further preimages over a sufficiently
small common target neighborhood.  Since the target plane is
simply connected and the source plane connected, path lifting
gives a single sheet.  The inverse is smooth by the local inverse
theorem.

Replace $f$ by $(f-f(0))/(f(1)-f(0))$.
It now fixes zero and one, retains its equation, and has derivative
tending to a nonzero complex scalar.  There is at most one such
normalized solution for a fixed $\nu$: if $f,g$ are two, the chain
rule shows that $g\circ f^{-1}$ is holomorphic.  Its derivative
tends at infinity to the quotient of the two limiting scalars,
so is a bounded entire function.  Liouville's theorem makes the
transition affine, and the two fixed points make it the identity.

Apply this uniqueness to $\widehat f=J\circ f\circ J$.
Away from zero the chain rule and the reflection identity for
$\nu$ show that $\widehat f$ solves the same equation.
The solution $f$ and its inverse are holomorphic off a compact
set, because their Beltrami coefficients vanish there.
Thus $\widehat f$ and its inverse are holomorphic in punctured
neighborhoods of zero.  Properness makes their values tend to
zero, so the removable singularity theorem extends both there;
their inverse identities persist.  They are therefore smooth
diffeomorphisms at zero as well.  Near zero, $f$ is holomorphic
with $f'(0)\ne0$.  As $|z|\to\infty$, differentiation of the
reflection gives the complex scalar
\[
 \widehat f'(z)=
 \frac{\overline{f'(Jz)}}
      {\overline{\,f(Jz)/(Jz)\,}^{\,2}}
 \longrightarrow\frac1{\overline{f'(0)}}\ne0.
\]
The reflected map fixes zero and one, so uniqueness proves
$fJ=Jf$.  The unit circle, the fixed set of $J$, is consequently
preserved.  No interior point can map outside it: the image of
a radial path from zero would first hit the circle, contradicting
injectivity and circle preservation.  Applying the same argument
to the inverse proves preservation of both the open and closed
disks.

Set $\Phi=f^{-1}$.  On the closed disk $\nu=\mu$, and the identity
$Df_{\Phi(z)}D\Phi_z=I$ gives
\[
 Q_{\Phi(z)}(D\Phi_zv,D\Phi_zv)
   =\frac{d(\Phi(z))}{|\partial f(\Phi(z))|^2}|v|^2.
\]
This is the promised coordinate formula.  Smoothness on the
whole plane bounds both first derivatives on the closed disk,
so both restrictions are Lipschitz.  The same construction
applies directly to any smooth Beltrami coefficient supported
away from the center and boundary, as required by the inner
variations below.

The use of smooth coefficients agrees with the scope of
\cite[Theorem 2.2, p.~6]{FitziWenger2019}.  Its epsilon-conformal
disk result is Theorem 4.1, pp.~8--10.  The present disk
preservation argument uses reflection and normalized uniqueness;
the cited proof instead conformally maps the image domain back
to the disk.  The matrix and energy computations here concern
Riemannian targets and half Dirichlet energy.

\paragraph{Area-to-energy comparison.}
The coordinates just constructed reparameterize $D$ to the
Lipschitz disk $D'=D\circ\Phi$ with the same area.  The chain rule,
the isothermal identity and change of variables give
\[
 E(D')=\int_{\mathbb D}\frac12\sqrt{\det Q}\,
                         \operatorname{tr}(Q^{-1}H).
\]
In this calculation $D\Phi(D\Phi)^T=cQ^{-1}$ and
$|\det D\Phi|=c/\sqrt{\det Q}$ cancel the conformal factor.
A conformal disk
automorphism puts the three labels at the chosen pins without
changing energy or area.  The two possible positions $\pm i$ retain
both boundary orientations.

Every such $D'$ belongs to the weak class with its actual weak
derivatives and Green trace.  Minimality gives
$E(u)\le\operatorname{Area}(D)+\eta$.  First let $\eta\downarrow0$,
then take the area infimum over $D$.  The same conformal normalization
of arbitrary weak disks shows that this normalized minimum minimizes
energy in the unpinned weak class as well.
\lean{PoincareMT.m65SpanningDisk_exists_normalized_energy_competitor}
\lean{PoincareMT.m65Plateau_weak_minimum}

We next obtain conformality in the original weak derivatives.  Put
$a=G_u(d_1u,d_1u)$, $b=G_u(d_1u,d_2u)$, $c=G_u(d_2u,d_2u)$,
using embedded derivatives in this formula, and set $e_0=(a+c)/2$.
For a smooth complex test function $\nu$ supported away from the
center and boundary, Beltrami reparameterization produces admissible
weak competitors with energies
\[
 V(t)=\int_{\mathbb D}
 \frac{(1+t^2|\nu|^2)e_0
       -t\{\operatorname{Re}\nu\,(a-c)+2\operatorname{Im}\nu\,b\}}
      {1-t^2|\nu|^2}.
\]
For sufficiently small $|t|$ the denominator is bounded below, and
the derivative is dominated by a constant times
$|e_0|+|a-c|+|b|$.  Minimality at $t=0$ therefore gives
\[
 \int_{\mathbb D}
       \{\operatorname{Re}\nu\,(a-c)+2\operatorname{Im}\nu\,b\}=0.
\]
Real and purely imaginary tests imply $a=c$ and $b=0$ almost
everywhere.  The omitted center and boundary have area measure zero.
No differentiable dependence of the chosen isothermal maps on $t$
is needed: their exact energy formula is what is differentiated.
\lean{PoincareMT.m65WeakDisk_conformal_of_normalized_minimum}

\paragraph{Regularity of the same minimum.}
Trace-preserving replacement on an interior disk proves local energy
minimality.  For almost every radius the actual weak map has an
absolutely continuous circle trace and the local Green identity.
If that trace has small oscillation, it lies in one uniformly
controlled target chart.  Coning the chart trace to its value at
one point gives a weak disk with the same trace and energy bounded
by a fixed multiple of its angular energy.  Gluing it into $u$ and
using minimality gives $E(r)\le C rE'(r)$.  If the trace oscillation
is too large for this construction, Cauchy--Schwarz gives a lower
bound for $rE'(r)$; the fixed outer energy bound yields the same
inequality after enlarging $C$.  These statements hold at the same
almost every radius, where the trace, Green, and radial energy
identities all hold.  Absolute continuity of $E$ gives
\[
 E(r)\le E(R)(r/R)^{1/C}.
\]
The constants can be fixed for centers in a smaller disk.  Averaging
and Poincare estimates then construct a Holder continuous
representative of the original weak map; continuity and closedness
of $e(M)$ keep its values in the target.

Supported variations in a captured target chart give the weak
harmonic-map equation $\Delta X=f(X,DX)$ with
$|f|\le C|DX|^2$.  Here the derivatives remain those of the original
weak map.  We describe the regularity step because the initial
$W^{1,2}$ bound alone does not supply continuous derivatives.
Write $e_X=|DX|^2$ and let $\beta>0$ be a Holder exponent for $X$,
reduced so that the uniform disk energy bound is $Cr^{2\beta}$.
Annular summation gives, uniformly for nearby centers $x$,
\[
 \int |z-x|^{-\beta}e_X(z)\,dz< C_\beta.
\]
On a sufficiently small neighborhood, continuity makes
$|f(z)\cdot(X(z)-X(x))|\le e_X(z)/2$.
Suppose the weighted energy bound holds with exponent
$\alpha\le1$.  The logarithmic potential for the weak Poisson
equation bounds the weighted first derivative with exponent
$s+1$, where $s=\alpha-\beta/2<\alpha$.  Test the equation
against the centered map times a cutoff and the regularized weight
$|z-x|^{-b}$, where $b=\alpha+\beta/2$.  The forcing pairing is
absorbed by the preceding half-energy bound.  Differentiating the
weight leaves a first derivative with weight $|z-x|^{-b-1}$;
the Holder bound for the centered map reduces this exponent to
$b+1-\beta=s+1$, just controlled by the potential estimate.
The cutoff terms have a uniform bound on a smaller disk.  This
proves the weighted energy bound with exponent $b$.
Starting at $\alpha=\beta$,
finitely many steps reach an exponent greater than one.
The gradient kernel of the logarithmic potential has size
$C|z-x|^{-1}$; the stronger weighted bound makes its near-diagonal
integral uniformly small.  Its far part is continuous by dominated
convergence.  The harmonic remainder is smooth, so the actual weak
derivatives have continuous representatives and $X$ is $C^1$.
The local harmonic-map bootstrap, the $\alpha=1$ case of the
regularity argument in Chapter~\ref{ch:filling-width-v4}, now makes
the same map smooth.  Its coordinate equation is the vanishing of
the geometric tension field.
\lean{PoincareMT.M65Interior.localMinimum_energy_power_decay}
\lean{PoincareMT.M65Boundary.quadratic_holder_contDiffAt}
\lean{PoincareMT.m65PlateauInteriorRegularityInput_proved}

At a boundary point use the holomorphic coordinate
$P_p(z)=p\exp(iz)$, $|p|=1$, from a small upper half-disk to the
original disk.  The boundary values lie in a regular target arc.
Here a cone to an arbitrary trace value would change the diameter
trace incorrectly.  Instead, straighten the captured target arc
to an axis and write $v(\theta)$, $0\le\theta\le\pi$, for the
absolutely continuous trace on a semicircle of radius $r$.
Its endpoint values lie on that axis.  With
$m=(v(0)+v(\pi))/2$, use the half-cone
\[
 v_r(s,\theta)=m+\frac{s}{r}\bigl(v(\theta)-m\bigr).
\]
Its diameter trace is the single affine interval from $v(\pi)$
to $v(0)$.  The boundary label is replaced on precisely this
source arc by the corresponding monotone target-arc label, and
is unchanged elsewhere.  Monotonicity at both endpoints shows
that the new label remains weakly monotone.  The actual traces
agree on the semicircle, so the two Green identities glue this
half-cone to the original map as an admissible weak disk.
Unpinned minimality and cancellation of the unchanged exterior
energy give $E^+(r)\le C r(E^+)'(r)$.  If the semicircle has
oscillation too large for one target chart, its angular energy
has a fixed positive lower bound and the outer energy bound
gives the same inequality.  The constants are uniform along
the compact boundary curve.

Integrating this inequality gives power decay on half-disks
centered on the diameter.  Even reflection constructs a weak
map on a full disk: the semicircle Green identities at common
good radii give its reflected weak derivatives.  A disk near
the diameter is bounded by a larger boundary-centered
half-disk, while disks away from the diameter use the interior
decay already proved.  Thus the reflected map has uniform
power decay on all sufficiently small disks.  No local
minimality of the reflected map is asserted.
The averaging construction yields a continuous
representative whose diameter trace is exactly
$\gamma\circ\beta\circ P_p$.  In the open half-disk it agrees
with the smooth interior representative: two continuous functions
equal almost everywhere are equal there.

Straighten the regular target arc to one coordinate axis.  Write
the resulting map as $X=(x,Y)$, with $Y=0$ on the diameter.
The coordinate metric remains smooth and positive definite.
Harmonicity and the two conformality equations give
$|\Delta X|\le C|DY|^2$.  Odd reflection of $Y$ has a genuine weak
Poisson equation across the diameter; the zero trace cancels the
boundary terms.  Its forcing is bounded by $C|DY|^2$, and its
Holder and energy-decay bounds follow from the original map.
The preceding weighted argument makes $DY$ continuous up to the
diameter.

Conformality recovers the remaining derivatives without dividing by
the differential of $X$.  In complex notation it is one quadratic
equation in the missing longitudinal column.  Completing the square
expresses its square as a continuous function of the already
controlled transverse columns and metric coefficients.  At a zero
of that function the modulus of the square root tends to zero.
At a nonzero value, a local continuous square root exists and
connectedness of a smaller open half-disk fixes the sign of the
existing interior root.  Thus the full differential extends
continuously.  Integrating it along segments gives within-$C^1$
regularity of $X$, and inverse coordinates return that regularity
to the original map.  Overlapping extensions agree in the interior
and hence on their common closed-domain boundary.  This constructs
one within-$C^1$ map on the entire closed disk with the same weak
values, derivatives, and continuous boundary trace.
\lean{PoincareMT.M65Boundary.weakDisk_boundary_local_contMDiff}
\lean{PoincareMT.m65PlateauBoundaryRegularityInput_proved}

\paragraph{Finite branches and a homeomorphic trace.}
In a harmonic target chart the complex differential
$\Phi=X_x-iX_y$ satisfies $\bar\partial\Phi=A\Phi$ with bounded
matrix coefficient locally.  On a small disk a Cauchy-transform
contraction solves $P=I+\mathcal C(AP)$ with $\|P-I\|<1$.
Consequently $P$ is invertible and $P^{-1}\Phi$ is holomorphic.
It either vanishes locally or factors as $z^m v(z)$ with $m\ge0$
and $v(0)\ne0$.

The boundary version uses the actual tangent of the regular Jordan
arc.  An invertible metric row frame makes the tangential component
of $\Phi$ real on the diameter and the transverse components purely
imaginary.  Reflection by the corresponding anti-complex involution
extends this framed field continuously to a full disk.  Its reflected
coefficient may jump across the diameter, but is bounded and smooth
on each open side.  The same small-disk Cauchy construction gives a
continuous invertible gauge.  After the gauge change the field is
holomorphic on both sides; rectangle integrals cancel along the
common continuous trace, proving holomorphicity across the diameter.
Undoing the frame gives, in the closed half-disk,
\begin{equation}
 X_x-iX_y=z^m Q(z),\qquad Q(0)\ne0.
 \label{eq:annular-boundary-factor}
\end{equation}
After shrinking, $Q$ is nowhere zero, continuous on the closed
half-disk, and $C^1$ on its interior.  The Cauchy-transform estimates
also give $DQ\in L^2$ and
$|Q(z)-Q(w)|\le C|z-w|^{1/2}$.  These estimates, rather than
continuity alone, will justify the curvature integral below.

Local vanishing of the differential would propagate through the
connected open disk by the interior alternative, making the map
constant.  This contradicts the surjective weakly monotone trace
onto a Jordan curve.  Thus every differential zero is isolated
relative to the closed disk.  The zero set is closed by within-$C^1$
regularity and compact, hence finite.  If $\beta$ were not injective,
weak monotonicity would give a nontrivial collapsed boundary arc.
The tangential derivative vanishes on that arc, and conformality,
extended to the boundary by continuity, makes the full differential
zero there.  Finiteness excludes this.  Therefore $\beta$ is a
continuous bijection of the circle, hence a homeomorphism.

Within-$C^1$ regularity on the compact disk supplies a Lipschitz
bound.  Area equals energy by conformality, and the same weak
derivative identities preserve that energy.  The resulting genuine
admissible disk $f$ consequently satisfies
\[
 A_g(\gamma)\le\operatorname{Area}_g(f)=E(u)\le A_g(\gamma),
\]
which finishes the attainment construction.
\lean{PoincareMT.M65StrictTrace.finite_disk_branches}
\lean{PoincareMT.m65Plateau_strict_trace}
\source{MorreyICM1950}\dcref{Sections 3--5, pp. 181--187}
\source{MT2007}\dcref{Lemma 19.2, pp. 437--439}

\subsection{Gauss--Bonnet with boundary branches}
\label{sec:annular-disk-gauss-bonnet}

Let $f$ be the disk just constructed and write
$f^*g=\lambda\,|dz|^2$ on its regular locus.  Its area density is
$\lambda$, which is continuous and nonnegative on the closed disk.
Define on the regular interior
\[
 G=-\tfrac12\Delta\log\lambda,
 \qquad L(R,\theta)=\tfrac12
       d\log\lambda_{Re^{i\theta}}(Re^{i\theta}).
\]
Let $W$ denote the curvature vector of the geometric boundary curve,
and define the continuous boundary flux
\[
 B(\theta)=\big\langle W(f(e^{i\theta})),
                   df_{e^{i\theta}}(e^{i\theta})\big\rangle_g.
\]
The differential here is the within differential on the closed
disk.  At a branch it is zero, so this expression remains defined.
On a regular boundary arc, $-B\,d\theta=k_g\,ds$, where $k_g$ is
the inward geodesic curvature.  We shall prove the form needed
without a boundary immersion assumption:
\begin{equation}
 2\pi\le\int_{\mathbb D}G-\int_{-\pi}^{\pi}B(\theta)\,d\theta.
 \label{eq:annular-branched-gauss}
\end{equation}

First, $G$ is absolutely integrable.  Near an interior branch the
complex differential has a finite power factor with a nonzero
$C^1$ residual.  Near a boundary branch use
(\ref{eq:annular-boundary-factor}).  Normal projection of the
covariant Hessian kills the derivative of $z^m$ because that term
is tangential.  Dividing its squared normal part by the conformal
factor therefore leaves a bound by
$C(1+|DQ|^2)$ on a smaller half-disk.  The residual is bounded
away from zero there.  The Gauss identity bounds the absolute
logarithmic curvature density by these integrable normal-Hessian
terms and a bounded ambient-curvature term times area density.
This proves local absolute integrability.  Under $P_p$, the local
factor is $(\lambda\circ P_p)|P'_p|^2$; since
$\Delta\log|P'_p|=0$, its logarithmic curvature density is
$(G\circ P_p)|P'_p|^2$.  The actual Jacobian change of variables
transports the local integrability to the original disk.  A finite
cover and the nullity of the finite branch set and boundary finish
the global assertion.

For the interior flux, put all interior branches $a_j$ inside some
$R_0<1$ and subtract their logarithmic singularities:
\[
 v=\tfrac12\log\lambda-\sum_jm_j\log|z-a_j|.
\]
The branch factors make $v$ $C^1$ across these centers.  Stokes'
formula, first off the centers and then using the integrable
curvature, gives
$\int_{|z|\le R}G=-R\int_{-\pi}^{\pi}dv_{Re^{i\theta}}(e^{i\theta})\,d\theta$
for $R_0<R<1$.  Each subtracted logarithm has nonnegative outward
radial derivative on such a circle, since
\[
 d\log|z-a_j|_z(z)=\frac{\langle z-a_j,z\rangle}{|z-a_j|^2}>0.
\]
Consequently
\begin{equation}
 -\int_{-\pi}^{\pi}L(R,\theta)\,d\theta
                 \le\int_{|z|\le R}G.
 \label{eq:annular-interior-flux}
\end{equation}

The boundary limit requires more than continuity of $Q$.  Its two
real residual columns are orthogonal and have the same positive
metric norm.  Normalize them to an orthonormal pair $(T,N)$ in the
tangent plane of the surface.  This pair extends continuously
through the branch, is Holder with exponent $1/2$, and has $L^2$
interior derivatives.  The branch order at a boundary point is
even: compose the boundary map with a local inverse parameter of
the Jordan arc.  The resulting real function is monotone, while
its derivative is $t^m q(t)$ with $q(0)\ne0$; an odd $m$ would
force a sign change.  Thus along the diameter the residual tangent
$T$ is, with a fixed sign, the unit tangent of the boundary curve.
It is $C^1$ there even at the branch.  Metric compatibility gives
for its boundary connection
\[
 J(t,0):=\langle\nabla_tT,N\rangle=-B(t+\arg p).
\]

At $z=t+ih$, $h>0$, differentiating the rotation from the residual
frame to the actual gradient frame gives
\[
 -\tfrac12\partial_h\log\lambda_{\rm loc}
       =J(t,h)+m\operatorname{Im}(1/z)\le J(t,h).
\]
Here $\operatorname{Im}(1/z)=-h/|z|^2\le0$.
The exponential coordinate satisfies
$P_p(t+ih)=e^{-h}e^{i(t+\arg p)}$ and $|P'_p|=e^{-h}$,
so the same identity reads
\begin{equation}
 L(e^{-h},t+\arg p)\le J(t,h)-1.
 \label{eq:annular-boundary-radial}
\end{equation}
This sign is the favorable boundary branch contribution.

Subdivide $[-\pi,\pi]$ into finitely many intervals lying in such
charts.  There is one common sequence $h_k\downarrow0$ on which
all their connection integrals converge to the diameter integrals.
Here is the estimate ensuring that assertion.  For each interval
replace the metric pairing with $N$ by its Euclidean metric dual
$V$; this dual is continuous and has $L^2$ tangential derivative.
Its slice energy $F(h)=\int|\partial_tV(t,h)|^2dt$ is integrable
in $h$.  For the finite sum of these energies choose heights with
$h_k\sum F(h_k)\to0$.  Otherwise this sum would eventually be
bounded below by a positive multiple of $1/h$, contradicting
integrability.  Integration by parts gives
\[
 \left|\int\langle\partial_t(T(\cdot,h)-T(\cdot,0)),
                                  V(\cdot,h)\rangle\right|
 \le C\sqrt h\left(1+\sqrt{\int|\partial_tV(t,h)|^2dt}\right).
\]
The endpoint term uses boundedness of $V$, and the integral term
uses the $1/2$-Holder bound for $T$ and Cauchy--Schwarz.  The
right side tends to zero at the chosen heights.  The remaining
fixed boundary derivative and the connection-coefficient terms
converge by continuity and compact domination.  Thus
$\sum\int J(t,h_k)\to-\int B$.

Integrate (\ref{eq:annular-boundary-radial}) over this subdivision
and combine with (\ref{eq:annular-interior-flux}):
\[
 2\pi-\sum\int J(t,h_k)
       \le-\int L(e^{-h_k},\theta)\,d\theta
       \le\int_{|z|\le e^{-h_k}}G.
\]
Absolute integrability of $G$ permits the last limit by dominated
convergence, proving (\ref{eq:annular-branched-gauss}).
\lean{PoincareMT.M65MinimalDisk.logarithmicGaussDensity_integrable}
\lean{PoincareMT.M65MinimalDisk.logarithmicGaussDensity_gaussBonnet}
\source{MT2007}\dcref{Lemma 19.2, pp. 438--439}

\subsection{The disk inequality for an immersed shortening curve}
\label{sec:immersed-disk-comparison}

The negative constant in the profile comes from the disk, not from the
annulus.  Let $c(x,t)$ be a smooth family of immersed null loops on an
open time interval compactly contained in $(a,b)$.  Suppose first that
$c(\cdot,t)$ is embedded at the time in question and that its velocity
$V$ satisfies $|V-H|\le\epsilon$.  An area-minimizing spanning disk is
conformal and smooth in the open disk. It is $C^1$ within the closed
disk, and its finite branch set may meet the boundary.
Fix the tested time $q$ and this attained disk $f$.  On a small
time interval the embedded regular loop remains embedded.
Local tubular charts and a partition of unity extend its actual
time velocity to a smooth time-dependent ambient field.
Compactness gives a common interval for its flow $\Psi_t$,
normalized by $\Psi_q=\mathrm{id}$.  Uniqueness for the ordinary
differential equation gives
$\Psi_t(c(x,q))=c(x,t)$ with the original parameter $x$.
Thus $f_t=\Psi_t\circ f$ is an actual Lipschitz spanning disk
at time $t$, with the same boundary homeomorphism as $f$.

Use the half-energy
\[
 E(t)=\frac12\int_D
    \bigl(|(f_t)_x|_{g(t)}^2+|(f_t)_y|_{g(t)}^2\bigr).
\]
The Gram determinant inequality gives
$A_{g(t)}(c(t))\le\operatorname{Area}_{g(t)}(f_t)\le E(t)$.
At $q$, conformality and attainment give
$E(q)=A_{g(q)}(c(q))$, including at the zero-rank points.
The ambient motion and the actual within-$C^1$ disk columns
make the energy and its time derivative continuous on the
compact disk-time product.  Differentiation under the integral
is therefore justified.  The metric term is minus the Ricci
trace on those columns.  The map variation is integrated by
parts using harmonicity: first on concentric smaller disks,
then to the boundary using the continuous first-jet trace.
The resulting outward radial flux is the pairing of the
ambient velocity with $f_r$ on $\partial D$.
This constructs a differentiable upper bound with equality at
$q$, so its derivative bounds the upper forward quotient of
the filling infimum there.

At a regular point choose an orthonormal tangent
pair $e_1,e_2$ and a normal $\nu$.  In dimension three,
\[
 \operatorname{Ric}(e_1,e_1)+\operatorname{Ric}(e_2,e_2)
     =R/2+\sec(e_1,e_2),\qquad
 K_\Sigma=\sec(e_1,e_2)-|\mathrm{II}|^2/2.
\]
The inward conormal component of $H$ is $k_g$.  Metric variation,
minimality, and the residual bound give
\begin{align*}
 D^+A_{g(t)}(c(t))
 &\le-\int_\Sigma(R/2+K_\Sigma+|\mathrm{II}|^2/2)
       -\int_{\partial\Sigma}k_g+\epsilon\operatorname{Length}(c(t))\\
 &\le-2\pi-\tfrac12R_{\min}(t)A_{g(t)}(c(t))
       +\epsilon\operatorname{Length}(c(t)).
\end{align*}
For a branched disk the integral notation in this calculation is
interpreted on the regular locus.  Precisely, let $Q$ be the Ricci
trace on its two differential columns.  It is integrable, and the
Gauss identity gives almost everywhere
$G\le Q-(R\circ f)\lambda/2$.
The metric and boundary first variation is bounded above by
$-\int Q+\int B+\epsilon\operatorname{Length}(c(t))$.
For this last error term, conformality on the closed trace
gives $|f_r|=|f_\theta|$ on the unit circle, including its
branch points.  Integration of the latter norm is the length
of the original embedded loop, because its boundary label is
a homeomorphism.  Pairing $V-H$ with the radial derivative
therefore contributes at most $\epsilon$ times that length.
Equation~(\ref{eq:annular-branched-gauss}) therefore gives the
displayed area inequality, including at disks with boundary
branches.  This uses the disk from Theorem~\ref{thm:annular-plateau}
in exactly the admissible filling class defining $A_g$.
\lean{PoincareMT.m65EmbeddedFillingAreaInequality_proved}
\source{MT2007}\dcref{Lemma 19.2 and Claim 19.5, pp. 437--441}

For immersed curves, construct generic perturbations on a slightly larger
compact time interval.  Uniform immersion excludes coincidences between
sufficiently close angular parameters.  On the compact set of remaining
pairs, use finitely many coordinate perturbations, localized by smooth
cutoffs, whose parameter derivatives span the normal space of the diagonal
in $M\times M$.  More explicitly, near each original coincidence
with separated angles, vary three target coordinate directions
at one angle using a cutoff which vanishes near the other.
The three corresponding parameter derivatives of the coordinate
difference form an invertible block.  A finite cover supplies
one finite parameter family, and this block remains invertible
on a neighborhood of each covered coincidence.  Shrinking the
parameter ball retains immersion and captures every new
coincidence in those neighborhoods: otherwise a convergent
sequence of uncaptured coincidences would give an uncaptured
original one.

In real angular coordinates the difference equation has three
components and the variables $(x,y,t)$ have dimension three.
Solve for an invertible block of three perturbation parameters
by the implicit function theorem.  Each local zero set then
has dimension equal to the total parameter dimension.  Apply
Sard's theorem to its projection onto parameter space.
Outside a null set of parameters that projection has invertible
derivative at every zero, so each fixed-parameter coincidence
is isolated.  A finite cover of the compact zero set, using
a closed smaller parameter ball and the separated-angle
domain, gives one null exceptional parameter set.  Each
remaining parameter has a compact discrete coincidence set,
and hence finitely many coincidence times.
Thus each perturbed family $c_j$ is embedded at all but finitely many times.

Choose the perturbation parameters tending to zero sufficiently fast that
$|\partial_tc_j-H_j|\le1/(j+1)$ on the compact interval and the lengths
are bounded by one constant $L$.  These assertions follow from continuity
of the curvature and velocity expressions in the spatial second and time
first derivatives, and the positive lower speed of the original immersion.
The perturbation paths give small annuli to the original curves.  Gluing
actual disks through these annuli gives filled perturbations, and filling-
area continuity gives $f_j(t)\to f(t)$ at each endpoint, where
$f_j(t)=A_{g(t)}(c_j(t))$.

Set
\[
 w(t)=\exp\!\left(\int_a^t R_{\min}(v)/2\,dv\right),\qquad
 Y_f(t)=w(t)f(t)+2\pi\int_a^t w(v)\,dv.
\]
On every interval avoiding the finite exceptional set, the preceding
inequality says that
$Y_{f_j}(t)-(L/(j+1))\int_a^t w$ is nonincreasing.  Continuity extends
this assertion to the finitely many interval endpoints, and concatenation
gives it on the whole compact interval.  Passing to the actual endpoint
area limits yields $Y_f(t)\le Y_f(s)$ for $s\le t$.
Dividing by $w(t)>0$ proves
\[
 f(t)\le\Phi_{F,s}(f(s),t),\qquad
 \Phi_{F,s}(A,t)=\frac{w(s)A-2\pi\int_s^t w(v)\,dv}{w(t)}.
\]
This is the required immersed comparison.  Neither genericity nor a
residual estimate is an extra hypothesis on the original shortening curve.
\lean{PoincareMT.m65Exists_generic_filled_approximations}
\lean{PoincareMT.m65ImmersedFillingAreaComparison_of_diskComparison}
\source{MT2007}\dcref{Lemma 19.4 and Claim 19.5, pp. 439--441}

\subsection{Good times and a fixed delayed grid}
\label{sec:good-time-grid}

Fix $\eta>0$.  We first prove a statement for one member $z$, with a
circumference cutoff allowed to depend on $z$: either its product length
is less than $\eta/2$ throughout a fixed final interval $[T,b]$, or
\[
 f_z(b)\le\Phi_F(f_z(a),b)+\eta/2.
\]
The final time $T$ and all discretization parameters will precede $z$
and $\rho$.  This distinction is what makes the later finite-net step
possible.

For $0<\rho<1$, the ramp estimates applied to the chosen polygonal
approximation give
\[
 \int_a^b E_z(t)\,dt\le D,\qquad
 E_z(t)=\int_{c_z(t)}|H|^2\,ds,\qquad
 D=(\sup_z\operatorname{Length}_{g(a)}\Gamma(z)+1)e^{K_2(b-a)}.
\]
The bound uses the original family's length supremum.  At $a$ the actual
ramp is the canonical lift of the selected flattened polygon, whose length
is at most that original supremum plus $1$; the analytic estimates are
applied to those same solutions.

We first record an elementary scalar comparison, including its error
accounting.  Suppose $0\le f\le A_*$ and $f$ is $L_*$-Lipschitz on
$[a,b]$.  Compactness and continuity bound $w,w'$ and give constants
$G,U>0$, depending only on $A_*,L_*$ and the flow, such that
\[
 |Y_f(t)-Y_f(s)|\le G|t-s|,\qquad w(t)\le U.
\]
For a partition $a=t_0\le\cdots\le t_m=b$, mark some cells good.
Assume their restarted comparisons have error $e\ge0$:
$f(t_{i+1})\le\Phi_{F,t_i}(f(t_i),t_{i+1})+e$.
Multiplication by $w(t_{i+1})$ makes their weighted increments at most
$Ue$.  Every other increment is at most $G(t_{i+1}-t_i)$.
If the total length of the other cells is at most $d$, telescoping all
increments gives
\[
 f(b)\le\Phi_F(f(a),b)+\frac{Gd+Ume}{w(b)}.
\]
Choose $d>0$ with $Gd<w(b)\eta/4$.  Once a grid with $m$ cells has
been chosen, choose $e>0$ with $Ume<w(b)\eta/4$.
No unweighted difference of two profiles is being added across cells.
\lean{PoincareMT.m65FiniteProfileComparison}
\source{MT2007}\dcref{Claims 19.26--19.27, pp. 457--459}

Choose $B>1$ large, $T\in(a,b)$ close to $b$, and $h>0$ small enough
that
\[
 D/B+2h+(b-T)<d,\qquad e^{K_2(b-T)}<4/3.
\]
Here $h$ is the eventual grid step.  To arrange the curvature estimates,
let $c_*=\delta_0r_0^2>0$ be the local regularization constant from the
ramp chapter.  Choose $0<r<1$ with
\[
 r\le(\eta/2)e^{-K_2(b-a)},\quad rB\le c_*^2,\quad T+c_*r^2<b,
\]
then choose an integer $n$ and $h>0$ with
$a+(n+2)h=T$ and $4h\le c_*r^2$, making $h$ still smaller if needed.
These choices depend on the family bounds but not on a member or a
circumference.

If a member is not short throughout $[T,b]$, the exponential length
estimate gives length at least $(\eta/2)e^{-K_2(b-a)}$ at every earlier
time in $[a,T]$.  In a reference cell $(a+jh,a+(j+1)h)$ call the index
$j$ good when some interior time has $E_z\le B$.  If no such time exists,
the energy integral on that cell is at least $Bh$.  Hence the total length
of bad reference cells is at most $D/B$.
At a good reference time, Cauchy--Schwarz gives
$\Theta(I)^2\le |I|E_z\le rB$ on every arc of length at most $r$.
The local derivative estimates apply from that time forward, after a
positive delay.  They bound the squared $i$th curvature jet on
\[
 [a+jh+3h/2,\ a+jh+7h/2]
 \quad\text{by}\quad C_i/(h/2)^{i+1}.
\]
This enlarged interval lies strictly inside $(a,b)$ and contains the
comparison cell $[a+(j+2)h,a+(j+3)h]$ in its interior.  The two initial
cells and the final gap $[T,b]$ are not declared good.  Together with
shifted bad reference cells, their total length is at most
$D/B+2h+(b-T)<d$.  Thus the estimate retains the endpoint gaps explicitly.
\lean{PoincareMT.m65Family_exists_fixed_good_grid}

Fix a member and one enlarged cell $[r_-,r_+]$, and take
$\rho_k\downarrow0$.  We explain the required filled smooth limit.
The left endpoint $r_-=a+jh+3h/2$ lies in $[a,T]$, so the
earlier length floor applies there.  The right endpoint can exceed
$T$; its lower length bound must instead be propagated.  On this
whole cell let $B_i$ bound the squared $i$th curvature jet.
For the speed $v$ of a shrinking curve,
\[
 \partial_t\log v=-\operatorname{Ric}(S,S)-|H|^2,
 \qquad |\partial_t\log v|\le K_2+B_0.
\]
Integration preserves a positive length lower bound throughout
the enlarged cell, even beyond $T$.

For each $k$, reparameterize once at $r_-$ to constant speed
$L_k(r_-)/(2\pi)$.  This is a smooth increasing lift
$\phi_k(x+2\pi)=\phi_k(x)+2\pi$, with $\phi_k(0)=0$;
keep it fixed in time.  The speed equation and the upper and
lower initial lengths give constants $0<v_-\le v_k\le v_+$
on the whole cell, independent of $k$.  Since the relabeling
is fixed in time, the curve still satisfies the pure-normal
shortening equation and retains all intrinsic curvature jets.

Write $u_k$ for its vertical unit-tangent component and $\ell_*>0$
for the propagated length floor.  Degree one and positivity give
$\int u_k\,ds=\rho_k$, and the parallel vertical vector gives
$|\partial_su_k|\le |H_k|\le\sqrt{B_0}$.  There is a point
with $u_k\le\rho_k/L_k$.  Integrating
$|(u_k^2)'|\le2\sqrt{B_0}u_k$ around the circle from that point
yields the uniform bound
\[
 0\le u_k(x,t)^2\le (\rho_k/\ell_*)^2
                       +2\sqrt{B_0}\rho_k\longrightarrow0.
\]
The projected spatial derivative therefore has squared norm
$v_k^2(1-u_k^2)$, uniformly bounded away from zero for large $k$.

To pass to a smooth limit, intrinsic bounds must control the
actual coordinate derivatives.  Cover $M$ by finitely many
compact chart cores.  In a core write $q_k$ for the projected
coordinates and retain the state $(t,q_k,v_k,u_k)$.  Spatial
covariant differentiation of $S,H,\nabla_SH,\ldots$ expresses
their coordinate derivatives using the next intrinsic jet,
the speed and the chart Christoffel symbols.  The speed
equation is triangular for this purpose: after $m$ spatial
derivatives, its highest speed derivative solves a linear
inhomogeneous equation whose other terms use already bounded
lower speed derivatives and intrinsic jets.  Constant initial
speed makes all its positive-order initial derivatives zero.
Induction and Gronwall therefore bound every spatial derivative
of this state on the finite chart cores.

The joint bounds follow from an explicit finite-order equation.
Suppressing $k$, set
\[
 H_q=v^{-2}\bigl(q_{xx}+\Gamma(q)(q_x,q_x)\bigr)
                  -v_xv^{-3}q_x,
 \quad C=\operatorname{Ric}(q_x/v,q_x/v)
                  +|H_q|^2+(u_x/v)^2.
\]
Here metric, Ricci and Christoffel coefficients are evaluated
at $(t,q)$.  The equations are
\[
 q_t=H_q,\qquad v_t=-Cv,\qquad
 u_t=v^{-2}u_{xx}-v_xv^{-3}u_x+Cu.
\]
Their right sides are smooth functions of the state and its
first two spatial derivatives on the domain $v>0$.
At every finite order the bounds just obtained put these jets
in a compact subset of that domain.  Differentiating the
equations repeatedly bounds all mixed derivatives on smaller
time-chart rectangles.  The common lower speed bound prevents
any denominator from approaching zero.

Bounded speed and velocity first give one subsequence converging
locally uniformly as maps into the compact manifold.  Near a
point of this continuous limit, a smaller rectangle and the
tail of the sequence lie in one chart core.  The joint derivative
bounds and Arzela--Ascoli give smooth local subsequential limits.
Their positions equal the already chosen continuous limit;
integration of derivatives identifies every derivative limit
with its derivatives.  Applying the same argument to every
subsequence gives convergence of each chart derivative along
the original selected subsequence.  Thus the local limits agree
on overlaps and define one smooth periodic curve family $q$.
The slope limit is zero, so the limiting identity
$|q_x|^2=v^2$ gives immersion.  In the equation for $q_t$,
the limiting $v$ is exactly the base speed; hence $q_t$ is
the base curvature vector.  This proves smooth immersed
curve shortening on $(r_-,r_+)$, with convergence on compact
subintervals containing both comparison endpoints.
\lean{PoincareMT.m65FamilyCell_normalization_bounds}
\lean{PoincareMT.m65Slope_sq_le_circumference}
\lean{PoincareMT.m65Projected_exists_smooth_curveShortening_limit}

At a fixed interior time, the normalized projected loops converge
in $C^1$.  Their filling areas equal those of the original
projected loops: keep each disk map and its area integral, and
compose only its boundary label with the inverse circle
homeomorphism induced by $\phi_k$.  The inverse operation
identifies the whole admissible area ranges.  A short geodesic
interpolation annulus from a normalized loop to its limit has
area tending to zero by the two-column estimate.  Gluing one
actual approximating disk through such an annulus first supplies
a disk for the limit.  The two directional gluing inequalities
then show convergence of the filling infima.  Both comparison
endpoints use the same selected subsequence, each in its own
fixed metric.  The immersed comparison just proved thus applies
to this genuinely filled limit.
\lean{PoincareMT.m65FamilyCell_exists_filled_limit}
\lean{PoincareMT.m65FillingArea_tendsto_of_C1}
\lean{PoincareMT.m65FillingArea_relabel}

Consequently, for the fixed $e>0$, all sufficiently small circumferences
satisfy the restarted comparison with error $e$ on the smaller cell.
Indeed, otherwise choose failing $\rho_j<1/(j+1)$, take the filled smooth
limit, and pass both endpoint areas through the continuous profile.  The
limit would satisfy $f(t)\ge\Phi_{F,s}(f(s),t)+e$, a contradiction.
The cutoff asserts this implication whenever that cell has the required
geometric bounds, so the good set may still depend on $\rho$.

Take the minimum of these cutoffs over the finitely many cells.  Apply the
weighted scalar comparison to the actual good cells for each remaining
$\rho$.  There are $n+3$ cells including the two initial cells and the
final gap.  The selected $d,e$ give error less than $\eta/2$.  This proves
the stated per-member alternative with the same $T$ for every member.
\lean{PoincareMT.m65FamilyCell_exists_comparison_cutoff}
\lean{PoincareMT.m65Family_pointwise_terminal_alternative}
\source{MT2007}\dcref{Lemma 19.25, Claim 19.28 and Remark 19.29, pp. 455--461}

\subsection{One circumference and class preservation}
\label{sec:terminal-transfer}

Curve shortening in the product preserves the projected family homotopy
class: continuity in time gives a homotopy from the initial approximation to
each included slice, and the static flattening is homotopic to the raw family.
The base projection decreases length because the product tangent norm is the
sum of the base square and a nonnegative circle square.  Thus
\begin{equation}
 \operatorname{Length}_{g(t)}(\operatorname{pr}_1 c_z(t))
 \le \operatorname{Length}_{g(t)+d\theta^2}(c_z(t)).
 \label{eq:projection-length}
\end{equation}
\lean{PoincareMT.m65ProjectedFamily_homotopic}
\lean{PoincareMT.m65ProjectedFamilyLength_le}
\lean{PoincareMT.m65Projection_tangentNorm_le}

We now turn the per-member cutoffs into one circumference.  First choose
thresholds for transferring each of the two alternatives along an initial
annulus.  Let $z$ be a member, $z_*$ a proposed node, and $U$ an annulus between
their canonical initial ramps.  Annular evolution in base dimension three
and the two disk-gluing inequalities imply
\[
 |f_z(t)-f_{z_*}(t)|\le e^{5K_0(t-a)}\operatorname{Area}(U).
\]
To justify passage through the evolving infimum, take a competitor within
an arbitrary positive error of that infimum, project and glue it, and then
let the error tend to zero.  Attainment of the original annulus infimum is
unnecessary.

For the area branch, set $E=e^{5K_0(b-a)}$ and $Q=1/w(b)>0$.
The initial-value identity for the profile is
$\Phi_F(A,b)-\Phi_F(B,b)=Q(A-B)$.  If
$f_{z_*}(b)\le\Phi_F(f_{z_*}(a),b)+\eta/2$, then
\[
 f_z(b)\le\Phi_F(f_z(a),b)+\eta/2
                   +(E+Q)\operatorname{Area}(U).
\]
Thus the uniform choice $\mu_{\rm area}=\eta/[4(E+Q)]$ transfers
the half-error alternative at the node to the full-error alternative
at the member.
\lean{PoincareMT.m65AreaLoopTransfer}
\source{MT2007}\dcref{Lemma 19.30, p. 462}

For the short branch, suppose the node's product length is less than
$\eta/2$ at every time of $[T,b]$.  Set
\[
 B_1=\frac{|D|+1}{b-T},\qquad
 r_1=\min\{\eta/2,(1/400)^2/B_1\}.
\]
The energy bound forces a time $t\in(T,b)$ with $E_z(t)\le B_1$;
otherwise integration would exceed $D$.  At that time every member subarc
of length at most $r_1$ has total curvature at most $1/400<1/200$.
Let $\epsilon_1>0$ be the small-annulus length threshold at scale $r_1$,
and choose $\mu_{\rm short}=\epsilon_1/(2E)$.
If $\operatorname{Area}(U)<\mu_{\rm short}$, annular evolution gives
least annular area less than $\epsilon_1$ at $t$, hence an actual annulus
of area less than $\epsilon_1$ there.

If the projected member at $b$ had length at least $\eta$, its product
length at $b$ would be at least $\eta$.  Since
$e^{K_2(b-T)}<4/3$, its product length at $t$ would be at least
$3\eta/4\ge r_1$.  Small-annulus comparison, with this member as first
boundary, would force the node's length at $t$ to be at least
$(3/4)(3\eta/4)=9\eta/16>\eta/2$, a contradiction.  The projected member
therefore has terminal length less than $\eta$.
\lean{PoincareMT.m65ShortLoopTransfer}
\source{MT2007}\dcref{Claim 19.32, pp. 463--464}

Apply the finite-net construction to the \emph{approximated} family with
$\mu=\min\{\mu_{\rm area},\mu_{\rm short}\}$.
For each of its finitely many nodes take the per-member cutoff from the
fixed-grid argument.  Choose a positive circumference less than $1$, the
net cutoff, and every node cutoff.  For every member the net gives a node
and an actual initial annulus of area less than $\mu$.  If the node has
the area alternative, use the first transfer; if it is short on $[T,b]$,
use the second.  This proves the common terminal alternative.
The order of choices is
\[
 \eta\ ;\ T,\text{grid}\ ;\ \mu\ ;\ \text{finite nodes}\ ;\
 \text{node cutoffs}\ ;\ \rho\ ;\ z\ ;\ \text{its node}.
\]
\lean{PoincareMT.m65CommonTerminalAlternative}
\lean{PoincareMT.m65CommonTerminalAlternative_proved}
\lean{PoincareMT.m65DeformedFamilyOfProjectedEstimate}
\source{MT2007}\dcref{Claims 19.25--19.32, pp. 453--464}

The deformation is the projected solution at this single circumference.
It is jointly continuous in time and the source parameter by continuous
dependence in the ramp construction.  The concatenation of the initial
approximation homotopy with the time homotopy
$s\mapsto\operatorname{pr}_1c_z(a+s(t-a))$ identifies every slice with
the original free family.  A loop homotopic to a null loop is null, so the
pointwise nullity follows at every time.  Initial filling-area accuracy is
the static approximation estimate, and choosing $\eta=\zeta$ gives
(\ref{eq:terminal-alternative}).

When $t_0=t_1$, the original family itself is the deformation.  Its initial
area error is zero and the initial-value identity for $\Phi_F$ supplies the
area branch.  When $t_0<t_1$, the preceding construction applies to the same
selected approximation and its curve estimates.  No connectedness,
vanishing of $\pi_2(M)$, or nonzero class hypothesis is required for this
deformation theorem.
\lean{PoincareMT.m65ZeroDurationDeformedFamily}
\lean{PoincareMT.m65Construction}
\lean{PoincareMT.m65LoopFamilyDeformation}
\lean{PoincareMT.m65LoopFamilyDeformation_from_predecessors}

\section{Passage to a fixed width class}
\label{sec:annular-interfaces}

Write $W_g(\Gamma)=\sup_{z\in S^2}A_g(\Gamma(z))$ and let
$W_g[\Gamma]$ be the infimum of these suprema over continuous pointwise-null
families freely homotopic to $\Gamma$.  Compactness of the parameter sphere
and continuity of filling area make each supremum finite; the competitor
class is nonempty because it contains $\Gamma$ itself.  Nonnegativity
therefore bounds its infimum below.

For any $\epsilon>0$ choose a class competitor $\Gamma^\epsilon$ with
$W_{g(a)}(\Gamma^\epsilon)<W_{g(a)}[\Gamma]+\epsilon$ and apply the
deformation to it.  Free homotopy preservation makes its terminal slice
an admissible competitor in the \emph{same} class at time $b$.
The initial-area error and positivity of the profile's initial-value
coefficient imply, on its area branch,
\[
 A_{g(b)}(\Gamma_b^\epsilon(z))
 \le\Phi_F(W_{g(a)}[\Gamma]+\epsilon+\zeta,b)+\zeta.
\]
The other branch has length less than $\zeta$.  The short-loop filling
estimate from the filling-width chapter supplies, for the fixed terminal
metric, a modulus $\omega_b(\zeta)\to0$ bounding its area.  Taking the
supremum and then the class infimum gives
\[
 W_{g(b)}[\Gamma]\le
 \max\{\omega_b(\zeta),
       \Phi_F(W_{g(a)}[\Gamma]+\epsilon+\zeta,b)+\zeta\}.
\]
This is the exact use of an initial near-minimizer: the deformation need
not begin at the original representative.  Based class labels, positivity
for a nonzero class, and the scalar-clock conversion used for extinction
are additional inputs of the next chapter.  In particular the maximum
above cannot be discarded merely from the terminal alternative.
